\documentclass[11pt]{article}

\usepackage[margin=1in]{geometry}
\usepackage{amsmath,amssymb,amsthm}
\usepackage{enumitem}
\usepackage{xcolor}
\usepackage{hyperref}

\hypersetup{
  colorlinks=true,
  linkcolor=blue!60!black,
  urlcolor=blue!60!black
}

\newcommand{\Prob}{\mathbb{P}}
\newcommand{\E}{\mathbb{E}}
\newcommand{\N}{\mathbb{N}}
\newcommand{\1}{\mathbf{1}}
\newcommand{\Knowledge}[1]{%
  \par\smallskip\noindent
  \textbf{Knowledge used.}
  {\small\color{blue!60!black}\raggedright #1\par}
  \par\smallskip
}

\title{Chapter 5 Exercises}
\author{\textsc{Solutions to Rick Durrett's}\\
\textit{Probability: Theory and Examples}}
\date{\today}

\begin{document}
\maketitle

\noindent
These solutions use the local Chapter 5 LLM/Viki knowledge set in
\path{Probability/LLM_Viki_Dataset/Chapter5_Knowledge}.  The cited
knowledge IDs refer to entries in
\path{chapter5_knowledge_pieces.jsonl}.

\section*{Exercise 5.1.1}

\noindent\textbf{Question.}
Let $\xi_1,\xi_2,\ldots$ be i.i.d. random variables taking values in
$\{1,2,\ldots,N\}$, each with probability $1/N$.  Show that
\[
  X_n=\left|\{\xi_1,\ldots,\xi_n\}\right|
\]
is a Markov chain and compute its transition probability.

\Knowledge{
\path{durrett_ch5_markov_property_countable};
\path{durrett_ch5_general_transition_probability}.
}

\noindent\textbf{Solution.}
The state $X_n$ is the number of distinct values seen among the first
$n$ observations.  If $X_n=k$, then exactly $k$ of the $N$ possible
values have already appeared.  The next observation $\xi_{n+1}$ is
independent of the past, so
\[
  \Prob(X_{n+1}=k+1\mid X_0,\ldots,X_n)
  =\frac{N-k}{N},
\]
because there are $N-k$ unseen values.  Similarly,
\[
  \Prob(X_{n+1}=k\mid X_0,\ldots,X_n)
  =\frac{k}{N},
\]
because drawing one of the already observed values does not change the
number of distinct values.  These probabilities depend on the past only
through $X_n$, so $(X_n)$ is a Markov chain.

If we include the natural initial state $X_0=0$, the state space is
$\{0,1,\ldots,N\}$ and the transition probability is
\[
  p(k,j)=
  \begin{cases}
    k/N, & j=k,\\
    (N-k)/N, & j=k+1,\\
    0, & \text{otherwise},
  \end{cases}
  \qquad 0\le k\le N,
\]
with the convention that $p(N,N)=1$.

\section*{Exercise 5.1.2}

\noindent\textbf{Question.}
Let $\xi_1,\xi_2,\ldots$ be i.i.d. random variables taking values in
$\{-1,1\}$, each with probability $1/2$.  Let $S_0=0$,
$S_n=\xi_1+\cdots+\xi_n$, and
\[
  X_n=\max\{S_m:0\le m\le n\}.
\]
Show that $X_n$ is not a Markov chain.

\Knowledge{
\path{durrett_ch5_markov_property_countable};
\path{durrett_ch5_random_walk_chain}.
}

\noindent\textbf{Solution.}
The record maximum $X_n$ does not remember the current location $S_n$.
That missing information affects whether the next step can create a new
maximum.

Consider the process through time $3$.  On the event
\[
  A=\{X_0=0,\ X_1=0,\ X_2=0,\ X_3=1\},
\]
we must have the increment sequence $(-1,+1,+1)$, so $S_3=1$ and
$X_3=1$.  Therefore
\[
  \Prob(X_4=2\mid A)=\Prob(\xi_4=1)=\frac12.
\]

Now consider the event
\[
  B=\{X_0=0,\ X_1=1,\ X_2=1,\ X_3=1\}.
\]
Then $\xi_1=1$, $\xi_2=-1$, and $\xi_3$ may be either $1$ or $-1$.
Conditional on $B$, these two possibilities are equally likely.  If
$\xi_3=1$, then $S_3=1$ and the next maximum reaches $2$ with
probability $1/2$.  If $\xi_3=-1$, then $S_3=-1$ and one step cannot
reach level $2$.  Hence
\[
  \Prob(X_4=2\mid B)
  =\frac12\cdot\frac12+\frac12\cdot 0
  =\frac14.
\]

Both events have $X_3=1$, but the conditional probabilities of
$X_4=2$ are different.  Thus the conditional law of the future depends
on more than the current value of $X_n$, and $(X_n)$ is not a Markov
chain.

\section*{Exercise 5.1.3}

\noindent\textbf{Question.}
Let $\xi_0,\xi_1,\ldots$ be i.i.d. random variables taking values in
$\{H,T\}$, each with probability $1/2$.  Show that
$X_n=(\xi_n,\xi_{n+1})$ is a Markov chain and compute its transition
probability $p$.  What is $p^2$?

\Knowledge{
\path{durrett_ch5_markov_property_countable};
\path{durrett_ch5_chapman_kolmogorov}.
}

\noindent\textbf{Solution.}
The state space is
\[
  \mathcal{S}=\{HH,HT,TH,TT\}.
\]
If $X_n=(a,b)$, then
\[
  X_{n+1}=(\xi_{n+1},\xi_{n+2})=(b,\xi_{n+2}).
\]
The variable $\xi_{n+2}$ is independent of the past and is equally
likely to be $H$ or $T$.  Therefore
\[
  p\bigl((a,b),(c,d)\bigr)
  =
  \begin{cases}
    1/2, & c=b,\\
    0, & c\ne b.
  \end{cases}
\]
Equivalently, from $(a,b)$ the chain moves to $(b,H)$ and $(b,T)$ with
probability $1/2$ each.

After two steps,
\[
  X_{n+2}=(\xi_{n+2},\xi_{n+3}),
\]
and $\xi_{n+2}$ and $\xi_{n+3}$ are independent fair coin tosses,
independent of $X_n$.  Hence
\[
  p^2\bigl((a,b),(c,d)\bigr)=\frac14
  \qquad
  \text{for all } (a,b),(c,d)\in\mathcal{S}.
\]
In the order $(HH,HT,TH,TT)$,
\[
  p=
  \begin{pmatrix}
  1/2 & 1/2 & 0 & 0\\
  0 & 0 & 1/2 & 1/2\\
  1/2 & 1/2 & 0 & 0\\
  0 & 0 & 1/2 & 1/2
  \end{pmatrix},
  \qquad
  p^2=
  \begin{pmatrix}
  1/4 & 1/4 & 1/4 & 1/4\\
  1/4 & 1/4 & 1/4 & 1/4\\
  1/4 & 1/4 & 1/4 & 1/4\\
  1/4 & 1/4 & 1/4 & 1/4
  \end{pmatrix}.
\]

\section*{Exercise 5.1.4}

\noindent\textbf{Question.}
Brother-sister mating.  In this scheme, two animals are mated, and among
their direct descendants two individuals of opposite sex are selected at
random.  These animals are mated and the process continues.  Suppose
each individual can be one of three genotypes $AA$, $Aa$, $aa$, and
suppose that the type of the offspring is determined by selecting a
letter from each parent.  With these rules, the pair of genotypes in the
$n$th generation is a Markov chain with six states:
\[
  AA,AA;\quad AA,Aa;\quad AA,aa;\quad Aa,Aa;\quad Aa,aa;\quad aa,aa.
\]
Compute its transition probability.

\Knowledge{
\path{durrett_ch5_markov_property_countable};
\path{durrett_ch5_general_transition_probability}.
}

\noindent\textbf{Solution.}
For a fixed parental pair, let
\[
  (p,q,r)
  =
  \bigl(\Prob(\text{offspring is }AA),
        \Prob(\text{offspring is }Aa),
        \Prob(\text{offspring is }aa)\bigr).
\]
The next brother-sister pair consists of two independent offspring,
recorded as an unordered genotype pair.  Therefore the transition
probabilities to the six states
\[
  AA,AA;\ AA,Aa;\ AA,aa;\ Aa,Aa;\ Aa,aa;\ aa,aa
\]
are
\[
  p^2,\quad 2pq,\quad 2pr,\quad q^2,\quad 2qr,\quad r^2.
\]

The Mendelian offspring distributions for the six parental states are
\[
\begin{array}{c|c}
\text{parental state} & (p,q,r)\\ \hline
AA,AA & (1,0,0)\\
AA,Aa & (1/2,1/2,0)\\
AA,aa & (0,1,0)\\
Aa,Aa & (1/4,1/2,1/4)\\
Aa,aa & (0,1/2,1/2)\\
aa,aa & (0,0,1).
\end{array}
\]
Thus, in the state order displayed above, the transition matrix is
\[
\begin{pmatrix}
1 & 0 & 0 & 0 & 0 & 0\\
1/4 & 1/2 & 0 & 1/4 & 0 & 0\\
0 & 0 & 0 & 1 & 0 & 0\\
1/16 & 1/4 & 1/8 & 1/4 & 1/4 & 1/16\\
0 & 0 & 0 & 1/4 & 1/2 & 1/4\\
0 & 0 & 0 & 0 & 0 & 1
\end{pmatrix}.
\]

\section*{Exercise 5.1.5}

\noindent\textbf{Question.}
Bernoulli-Laplace model of diffusion.  Suppose two urns, which we call
left and right, have $m$ balls each.  There are $b$ black balls, with
$b\le m$, and $2m-b$ white balls.  At each time, we pick one ball from
each urn and interchange them.  Let the state at time $n$ be the number
of black balls in the left urn.  Compute the transition probability.

\Knowledge{
\path{durrett_ch5_markov_property_countable};
\path{durrett_ch5_ehrenfest_chain};
\path{durrett_ch5_general_transition_probability}.
}

\noindent\textbf{Solution.}
Since $b\le m$, the possible states are $0,1,\ldots,b$.  If the current
state is $i$, then the left urn contains $i$ black and $m-i$ white
balls.  The right urn contains $b-i$ black and $m-b+i$ white balls.

The number of black balls in the left urn increases by $1$ exactly when
we choose a white ball from the left urn and a black ball from the right
urn.  Hence
\[
  p(i,i+1)
  =
  \frac{m-i}{m}\cdot\frac{b-i}{m}.
\]
It decreases by $1$ exactly when we choose a black ball from the left urn
and a white ball from the right urn.  Hence
\[
  p(i,i-1)
  =
  \frac{i}{m}\cdot\frac{m-b+i}{m}.
\]
It stays fixed when the two selected balls have the same color:
\[
  p(i,i)
  =
  \frac{i}{m}\cdot\frac{b-i}{m}
  +
  \frac{m-i}{m}\cdot\frac{m-b+i}{m}.
\]
All other transition probabilities are zero.  The boundary cases are
included automatically because the relevant factors vanish when
$i=0$ or $i=b$.

\section*{Exercise 5.1.6}

\noindent\textbf{Question.}
Let $\theta,U_1,U_2,\ldots$ be independent and uniform on $(0,1)$.  Let
$X_i=1$ if $U_i\le\theta$ and $X_i=-1$ if $U_i>\theta$, and let
$S_n=X_1+\cdots+X_n$.  In words, we first pick $\theta$ according to the
uniform distribution and then flip a coin with probability $\theta$ of
heads to generate a random walk.  Compute
\[
  \Prob(X_{n+1}=1\mid X_1,\ldots,X_n)
\]
and conclude that $S_n$ is a temporally inhomogeneous Markov chain.
This is due to the fact that ``$S_n$ is a sufficient statistic for
estimating $\theta$.''

\Knowledge{
\path{durrett_ch5_markov_property_countable};
\path{durrett_ch5_random_walk_chain}.
}

\noindent\textbf{Solution.}
Let
\[
  H_n=\#\{1\le i\le n:X_i=1\}.
\]
Then $H_n=(n+S_n)/2$.  Conditional on $\theta$, the variables
$X_1,\ldots,X_n$ are i.i.d. with
\[
  \Prob(X_i=1\mid\theta)=\theta.
\]
Since the prior distribution of $\theta$ is uniform on $(0,1)$, the
posterior density of $\theta$ given $X_1,\ldots,X_n$ is proportional to
\[
  \theta^{H_n}(1-\theta)^{n-H_n}.
\]
Thus the posterior distribution is
\[
  \theta\mid X_1,\ldots,X_n
  \sim \operatorname{Beta}(H_n+1,n-H_n+1).
\]
Therefore the predictive probability of a head on the next toss is the
posterior mean:
\[
  \Prob(X_{n+1}=1\mid X_1,\ldots,X_n)
  =
  \E(\theta\mid X_1,\ldots,X_n)
  =
  \frac{H_n+1}{n+2}.
\]
In terms of $S_n$,
\[
  \Prob(X_{n+1}=1\mid X_1,\ldots,X_n)
  =
  \frac{(n+S_n)/2+1}{n+2}
  =
  \frac{n+S_n+2}{2n+4}.
\]

Consequently,
\[
  \Prob(S_{n+1}=s+1\mid S_0,\ldots,S_n=s)
  =
  \frac{(n+s)/2+1}{n+2},
\]
and
\[
  \Prob(S_{n+1}=s-1\mid S_0,\ldots,S_n=s)
  =
  1-\frac{(n+s)/2+1}{n+2}
  =
  \frac{(n-s)/2+1}{n+2}.
\]
These transition probabilities depend only on the current value $s$ and
the time $n$, not on the earlier path.  Hence $(S_n)$ is a temporally
inhomogeneous Markov chain.

\section*{Exercise 5.2.1}

\noindent\textbf{Question.}
Let $A\in\sigma(X_0,\ldots,X_n)$ and
$B\in\sigma(X_n,X_{n+1},\ldots)$.  Use the Markov property to show that
for any initial distribution $\mu$,
\[
  \Prob_\mu(A\cap B\mid X_n)
  =
  \Prob_\mu(A\mid X_n)\Prob_\mu(B\mid X_n).
\]
In words, the past and future are conditionally independent given the
present.  Hint: Write the left-hand side as
$\E_\mu(\E_\mu(\1_A\1_B\mid \mathcal{F}_n)\mid X_n)$.

\Knowledge{
\path{durrett_ch5_markov_property_shift};
\path{durrett_ch5_general_transition_probability}.
}

\noindent\textbf{Solution.}
Since $A\in\mathcal{F}_n$, the tower property gives
\[
  \Prob_\mu(A\cap B\mid X_n)
  =
  \E_\mu\!\left(\E_\mu(\1_A\1_B\mid\mathcal{F}_n)\mid X_n\right)
  =
  \E_\mu\!\left(\1_A\E_\mu(\1_B\mid\mathcal{F}_n)\mid X_n\right).
\]
The event $B$ is a measurable function of the future path
$(X_n,X_{n+1},\ldots)$.  By the Markov property, there is a measurable
function $\varphi$ such that
\[
  \E_\mu(\1_B\mid\mathcal{F}_n)=\varphi(X_n).
\]
Therefore
\[
  \Prob_\mu(A\cap B\mid X_n)
  =
  \E_\mu(\1_A\varphi(X_n)\mid X_n)
  =
  \varphi(X_n)\E_\mu(\1_A\mid X_n).
\]
Also,
\[
  \Prob_\mu(B\mid X_n)
  =
  \E_\mu(\E_\mu(\1_B\mid\mathcal{F}_n)\mid X_n)
  =
  \varphi(X_n).
\]
Combining the last two displays gives
\[
  \Prob_\mu(A\cap B\mid X_n)
  =
  \Prob_\mu(A\mid X_n)\Prob_\mu(B\mid X_n).
\]

\section*{Exercise 5.2.2}

\noindent\textbf{Question.}
Let $X_n$ be a Markov chain.  Use Levy's $0$-$1$ law to show that if
\[
  \Prob\left(\bigcup_{m=n+1}^\infty \{X_m\in B_m\}\,\middle|\,X_n\right)
  \ge \delta>0
  \quad\text{on } \{X_n\in A_n\},
\]
then
\[
  \Prob(\{X_n\in A_n\ \text{i.o.}\}-\{X_n\in B_n\ \text{i.o.}\})=0.
\]

\Knowledge{
\path{durrett_ch5_markov_property_shift};
\path{durrett_ch5_hitting_return_decomposition}.
}

\noindent\textbf{Solution.}
For $N\ge0$, set
\[
  D_N=\bigcup_{m=N+1}^\infty \{X_m\in B_m\}.
\]
If $n>N$, then
\[
  D_N\supseteq \bigcup_{m=n+1}^\infty \{X_m\in B_m\}.
\]
Hence, using the Markov property and the hypothesis,
\[
  \Prob(D_N\mid\mathcal{F}_n)
  \ge
  \Prob\left(\bigcup_{m=n+1}^\infty \{X_m\in B_m\}\,\middle|\,\mathcal{F}_n\right)
  =
  \Prob\left(\bigcup_{m=n+1}^\infty \{X_m\in B_m\}\,\middle|\,X_n\right)
  \ge \delta
\]
on $\{X_n\in A_n\}$.

By Levy's $0$-$1$ law,
\[
  \Prob(D_N\mid\mathcal{F}_n)\longrightarrow \1_{D_N}
  \qquad\text{a.s.}
\]
If $\{X_n\in A_n\ \text{i.o.}\}$ occurs, then along infinitely many
$n>N$ the conditional probability above is at least $\delta$.  Therefore
$\1_{D_N}=1$ on this event.  Since this holds for every $N$,
\[
  \{X_n\in A_n\ \text{i.o.}\}
  \subseteq
  \bigcap_{N=0}^\infty D_N
  =
  \{X_n\in B_n\ \text{i.o.}\}
\]
up to a null set.  This proves the desired result.

\section*{Exercise 5.2.3}

\noindent\textbf{Question.}
A state $a$ is called absorbing if $P_a(X_1=a)=1$.  Let
\[
  D=\{X_n=a\ \text{for some }n\ge1\}
\]
and let $h(x)=P_x(D)$.  Use the result of the previous exercise to
conclude that $h(X_n)\to0$ a.s. on $D^c$.  Here a.s. means
$P_\mu$-a.s. for any initial distribution $\mu$.

\Knowledge{
\path{durrett_ch5_markov_property_countable};
\path{durrett_ch5_hitting_return_decomposition};
\path{durrett_ch5_strong_markov_property}.
}

\noindent\textbf{Solution.}
Fix $\varepsilon>0$ and define
\[
  A_n=\{x:h(x)\ge\varepsilon\},
  \qquad
  B_n=\{a\}.
\]
If $X_n=x$ and $h(x)\ge\varepsilon$, then the probability of hitting
$a$ at some future positive time is at least $\varepsilon$:
\[
  \Prob\left(\bigcup_{m=n+1}^\infty\{X_m=a\}\,\middle|\,X_n=x\right)
  =h(x)\ge\varepsilon.
\]
Since $a$ is absorbing, hitting $a$ once implies $X_m=a$ for infinitely
many $m$.  Exercise 5.2.2 gives
\[
  \Prob(\{h(X_n)\ge\varepsilon\ \text{i.o.}\}
  -\{X_n=a\ \text{i.o.}\})=0.
\]
On $D^c$, the chain never hits $a$ at any positive time, so
$\{X_n=a\ \text{i.o.}\}$ does not occur.  Hence
\[
  \Prob(\{h(X_n)\ge\varepsilon\ \text{i.o.}\}\cap D^c)=0.
\]
Taking $\varepsilon$ through the positive rationals shows that
$\limsup_n h(X_n)=0$ on $D^c$.  Since $h\ge0$, this is exactly
$h(X_n)\to0$ a.s. on $D^c$.

\section*{Exercise 5.2.4}

\noindent\textbf{Question.}
First entrance decomposition.  Let
$T_y=\inf\{n\ge1:X_n=y\}$.  Show that
\[
  p^n(x,y)
  =
  \sum_{m=1}^n P_x(T_y=m)p^{\,n-m}(y,y).
\]

\Knowledge{
\path{durrett_ch5_strong_markov_property};
\path{durrett_ch5_hitting_return_decomposition};
\path{durrett_ch5_chapman_kolmogorov}.
}

\noindent\textbf{Solution.}
On the event $\{X_n=y\}$, the first positive entrance time into $y$ must
be some $m\in\{1,\ldots,n\}$.  These events are disjoint, so
\[
  p^n(x,y)
  =
  P_x(X_n=y)
  =
  \sum_{m=1}^n P_x(T_y=m,\ X_n=y).
\]
By the strong Markov property at $T_y$, conditional on $T_y=m$ the chain
restarts from $y$ at time $m$.  Thus
\[
  P_x(X_n=y\mid T_y=m)=P_y(X_{n-m}=y)=p^{\,n-m}(y,y),
\]
and therefore
\[
  p^n(x,y)
  =
  \sum_{m=1}^n P_x(T_y=m)p^{\,n-m}(y,y).
\]

\section*{Exercise 5.2.5}

\noindent\textbf{Question.}
Show that
\[
  \sum_{m=0}^n P_x(X_m=x)
  \ge
  \sum_{m=k}^{n+k} P_x(X_m=x).
\]

\Knowledge{
\path{durrett_ch5_chapman_kolmogorov};
\path{durrett_ch5_hitting_return_decomposition}.
}

\noindent\textbf{Solution.}
For any state $y$, the expected number of visits to $x$ during the first
$n$ steps starting from $y$ is at most the corresponding quantity
starting from $x$.  Indeed, let $T_x=\inf\{m\ge0:X_m=x\}$.  If $y=x$,
there is equality.  If $y\ne x$, then by decomposing at the first visit
to $x$,
\[
  \sum_{r=0}^n P_y(X_r=x)
  =
  \sum_{\ell=1}^n P_y(T_x=\ell)
       \sum_{s=0}^{n-\ell} P_x(X_s=x)
  \le
  \sum_{s=0}^n P_x(X_s=x).
\]
Now use the Markov property at time $k$:
\[
\begin{aligned}
  \sum_{m=k}^{n+k}P_x(X_m=x)
  &=
  \sum_{r=0}^n P_x(X_{k+r}=x)\\
  &=
  \sum_y P_x(X_k=y)\sum_{r=0}^n P_y(X_r=x)\\
  &\le
  \sum_y P_x(X_k=y)\sum_{r=0}^n P_x(X_r=x)\\
  &=
  \sum_{r=0}^n P_x(X_r=x).
\end{aligned}
\]

\section*{Exercise 5.2.6}

\noindent\textbf{Question.}
Let $T_C=\inf\{n\ge1:X_n\in C\}$.  Suppose that $S-C$ is finite and for
each $x\in S-C$,
\[
  P_x(T_C<\infty)>0.
\]
Then there are $N<\infty$ and $\epsilon>0$ such that for all
$y\in S-C$,
\[
  P_y(T_C>kN)\le (1-\epsilon)^k.
\]

\Knowledge{
\path{durrett_ch5_strong_markov_property};
\path{durrett_ch5_hitting_return_decomposition}.
}

\noindent\textbf{Solution.}
For each $y\in S-C$, the assumption $P_y(T_C<\infty)>0$ implies that
there is an integer $N_y$ with
\[
  P_y(T_C\le N_y)>0.
\]
Since $S-C$ is finite, choose
\[
  N=\max_{y\in S-C}N_y,
  \qquad
  \epsilon=\min_{y\in S-C}P_y(T_C\le N).
\]
Then $N<\infty$ and $\epsilon>0$.

For $y\in S-C$,
\[
  P_y(T_C>N)\le1-\epsilon.
\]
Using the Markov property at times $N,2N,\ldots$, on the event
$\{T_C>jN\}$ the chain is still in $S-C$, so the conditional probability
of avoiding $C$ for the next $N$ steps is at most $1-\epsilon$.  Thus
\[
  P_y(T_C>(j+1)N\mid\mathcal{F}_{jN})
  \le 1-\epsilon
  \qquad\text{on }\{T_C>jN\}.
\]
Induction gives
\[
  P_y(T_C>kN)\le(1-\epsilon)^k.
\]

\section*{Exercise 5.2.7}

\noindent\textbf{Question.}
Exit distributions.  Let $V_C=\inf\{n\ge0:X_n\in C\}$ and let
\[
  h(x)=P_x(V_A<V_B).
\]
Suppose $A\cap B=\varnothing$, $S-(A\cup B)$ is finite, and
$P_x(V_{A\cup B}<\infty)>0$ for all $x\in S-(A\cup B)$.
\begin{enumerate}[label=(\roman*)]
\item Show that
\[
  h(x)=\sum_y p(x,y)h(y)\qquad x\notin A\cup B.
\]
\item Show that if $h$ satisfies this equation, then
$h(X_{n\wedge V_{A\cup B}})$ is a martingale.
\item Use this and Exercise 5.2.6 to conclude that
$h(x)=P_x(V_A<V_B)$ is the only solution of this equation that is $1$
on $A$ and $0$ on $B$.
\end{enumerate}

\Knowledge{
\path{durrett_ch5_expectation_operator_notation};
\path{durrett_ch5_hitting_return_decomposition};
\path{durrett_ch5_strong_markov_property}.
}

\noindent\textbf{Solution.}
Let $\tau=V_{A\cup B}$.  If $x\notin A\cup B$, then after one step the
chain is at $y$ with probability $p(x,y)$, and the future exit
probability is $h(y)$.  Therefore
\[
  h(x)=\sum_y p(x,y)h(y),
  \qquad x\notin A\cup B.
\]
Also $h=1$ on $A$ and $h=0$ on $B$.

Now suppose $h$ is any function satisfying the displayed equation off
$A\cup B$.  If $n\ge\tau$, then
$h(X_{(n+1)\wedge\tau})=h(X_{n\wedge\tau})$.  If $n<\tau$, then
$X_n\notin A\cup B$, so
\[
  \E\left(h(X_{(n+1)\wedge\tau})\mid\mathcal{F}_n\right)
  =
  \sum_y p(X_n,y)h(y)
  =
  h(X_n)
  =
  h(X_{n\wedge\tau}).
\]
Thus $h(X_{n\wedge\tau})$ is a martingale.

By Exercise 5.2.6, $\tau<\infty$ a.s. and in fact its tail decays
geometrically.  Since $S-(A\cup B)$ is finite and the boundary values
are fixed, any solution $h$ is bounded.  Taking expectations in the
stopped martingale identity gives
\[
  h(x)=E_x h(X_{n\wedge\tau}).
\]
Letting $n\to\infty$ and using bounded convergence,
\[
  h(x)=E_x h(X_\tau)
  =
  P_x(X_\tau\in A)
  =
  P_x(V_A<V_B).
\]
Therefore the exit probability is the unique solution with boundary
values $1$ on $A$ and $0$ on $B$.

\section*{Exercise 5.2.8}

\noindent\textbf{Question.}
Let $X_n$ be a Markov chain with $S=\{0,1,\ldots,N\}$ and suppose that
$X_n$ is a martingale.  Let $V_x=\min\{n\ge0:X_n=x\}$ and suppose
$P_x(V_0\wedge V_N<\infty)>0$ for all $x$.  Show that
\[
  P_x(V_N<V_0)=\frac{x}{N}.
\]

\Knowledge{
\path{durrett_ch5_expectation_operator_notation};
\path{durrett_ch5_hitting_return_decomposition}.
}

\noindent\textbf{Solution.}
Let
\[
  h(x)=P_x(V_N<V_0).
\]
By Exercise 5.2.7, $h$ is the unique function that is harmonic on
$\{1,\ldots,N-1\}$ and satisfies $h(0)=0$, $h(N)=1$.

Since $(X_n)$ is a martingale,
\[
  E_x X_1=x.
\]
Hence the function $f(x)=x/N$ satisfies
\[
  \sum_y p(x,y)f(y)
  =
  \frac1N\sum_y p(x,y)y
  =
  \frac{x}{N}
  =
  f(x),
  \qquad 1\le x\le N-1.
\]
It also has boundary values $f(0)=0$ and $f(N)=1$.  By uniqueness,
\[
  h(x)=f(x)=\frac{x}{N}.
\]

\section*{Exercise 5.2.9}

\noindent\textbf{Question.}
Wright-Fisher model.  Suppose $S=\{0,1,\ldots,N\}$ and consider
\[
  p(i,j)=
  \binom Nj
  \left(\frac{i}{N}\right)^j
  \left(1-\frac{i}{N}\right)^{N-j}.
\]
Show that this chain satisfies the hypotheses of Exercise 5.2.8.

\Knowledge{
\path{durrett_ch5_wright_fisher_chain};
\path{durrett_ch5_markov_property_countable}.
}

\noindent\textbf{Solution.}
Given $X_n=i$, the next state has distribution
\[
  X_{n+1}\sim \operatorname{Binomial}\left(N,\frac{i}{N}\right).
\]
Therefore
\[
  E(X_{n+1}\mid X_n=i)=N\cdot\frac{i}{N}=i.
\]
Thus $(X_n)$ is a martingale.

The state space is $\{0,1,\ldots,N\}$.  Also $0$ and $N$ are absorbing:
if $i=0$, then $X_{n+1}=0$ a.s.; if $i=N$, then $X_{n+1}=N$ a.s.  For
$1\le i\le N-1$,
\[
  p(i,0)=\left(1-\frac{i}{N}\right)^N>0,
  \qquad
  p(i,N)=\left(\frac{i}{N}\right)^N>0.
\]
Hence from every interior state there is positive probability of hitting
$0$ or $N$ in one step.  Thus
$P_i(V_0\wedge V_N<\infty)>0$ for all $i$, and all hypotheses of
Exercise 5.2.8 are satisfied.

\section*{Exercise 5.2.10}

\noindent\textbf{Question.}
In brother-sister mating described in Exercise 5.1.4, $AA,AA$ and
$aa,aa$ are absorbing states.  Show that the number of $A$'s in the pair
is a martingale and use this to compute the probability of getting
absorbed in $AA,AA$ starting from each of the states.

\Knowledge{
\path{durrett_ch5_markov_property_countable};
\path{durrett_ch5_hitting_return_decomposition};
\path{durrett_ch5_general_transition_probability}.
}

\noindent\textbf{Solution.}
Let $M_n$ be the number of $A$ alleles in the genotype pair at generation
$n$.  Thus $M_n$ takes values in $\{0,1,2,3,4\}$.  Given a parental pair
with $K$ total $A$ alleles, one offspring receives one randomly chosen
allele from each parent.  If the first parent has $k_1$ $A$ alleles and
the second has $k_2$ $A$ alleles, then $K=k_1+k_2$ and the expected
number of $A$ alleles in one offspring is
\[
  \frac{k_1}{2}+\frac{k_2}{2}=\frac{K}{2}.
\]
The next brother-sister pair consists of two offspring, so
\[
  E(M_{n+1}\mid\mathcal{F}_n)=2\cdot\frac{M_n}{2}=M_n.
\]
Thus $(M_n)$ is a martingale.

The only absorbing states are $AA,AA$ and $aa,aa$.  From each
nonabsorbing state there is positive probability of eventually reaching
one of them, and the state space is finite, so the absorption time is
finite a.s.  If
\[
  q=\Prob(\text{absorption in }AA,AA),
\]
then at absorption the martingale has value $4$ with probability $q$ and
$0$ with probability $1-q$.  Optional stopping gives
\[
  M_0=E(M_T)=4q,
  \qquad\text{so}\qquad q=\frac{M_0}{4}.
\]
The absorption probabilities in the state order
\[
  AA,AA;\quad AA,Aa;\quad AA,aa;\quad Aa,Aa;\quad Aa,aa;\quad aa,aa
\]
are therefore
\[
  1,\quad \frac34,\quad \frac12,\quad \frac12,\quad \frac14,\quad 0.
\]

\section*{Exercise 5.2.11}

\noindent\textbf{Question.}
Exit times.  Let $V_A=\inf\{n\ge0:X_n\in A\}$ and
$g(x)=E_xV_A$.  Suppose that $S-A$ is finite and for each $x\in S-A$,
$P_x(V_A<\infty)>0$.
\begin{enumerate}[label=(\roman*)]
\item Show that
\[
  g(x)=1+\sum_y p(x,y)g(y)\qquad x\notin A.
\]
\item Show that if $g$ satisfies this equation, then
$g(X_{n\wedge V_A})+n\wedge V_A$ is a martingale.
\item Use this to conclude that $g(x)=E_xV_A$ is the only solution of
this equation that is $0$ on $A$.
\end{enumerate}

\Knowledge{
\path{durrett_ch5_expectation_operator_notation};
\path{durrett_ch5_hitting_return_decomposition};
\path{durrett_ch5_strong_markov_property}.
}

\noindent\textbf{Solution.}
Let $\tau=V_A$.  Exercise 5.2.6 implies that $\tau<\infty$ a.s. and has
a geometrically decaying tail, so $E_x\tau<\infty$ for all $x$.

If $x\notin A$, then after one step the chain moves to $y$ with
probability $p(x,y)$.  One unit of time has elapsed, and the remaining
expected time is $g(y)$.  Thus
\[
  g(x)=1+\sum_y p(x,y)g(y),
  \qquad x\notin A.
\]
Also $g(x)=0$ for $x\in A$.

Now suppose $g$ satisfies this equation and is $0$ on $A$.  Let
\[
  M_n=g(X_{n\wedge\tau})+n\wedge\tau.
\]
If $n\ge\tau$, then $M_{n+1}=M_n$.  If $n<\tau$, then $X_n\notin A$ and
\[
\begin{aligned}
  E(M_{n+1}\mid\mathcal{F}_n)
  &=
  E(g(X_{n+1})+n+1\mid\mathcal{F}_n)\\
  &=
  \sum_y p(X_n,y)g(y)+n+1\\
  &=g(X_n)+n=M_n.
\end{aligned}
\]
Therefore $(M_n)$ is a martingale.

Taking expectations gives
\[
  g(x)=E_x\bigl[g(X_{n\wedge\tau})+n\wedge\tau\bigr].
\]
Letting $n\to\infty$, we have $X_\tau\in A$, so $g(X_\tau)=0$, and
$n\wedge\tau\to\tau$.  Hence
\[
  g(x)=E_x\tau=E_xV_A.
\]
Thus the expected hitting time is the unique solution with boundary
value $0$ on $A$.

\section*{Exercise 5.2.12}

\noindent\textbf{Question.}
In this problem we imagine that you are bored and start flipping a coin
to see how many tosses $N(HH)$ you need until you get two heads in a
row.  Let $\xi_0,\xi_1,\ldots$ be i.i.d. random variables taking values
in $\{H,T\}$, each with probability $1/2$, and let
$X_n=(\xi_n,\xi_{n+1})$ be the Markov chain from Exercise 5.1.3.  Use
the result of the previous exercise to compute
\[
  g(x)=\text{the expected time to reach state }(H,H)\text{ starting from }x.
\]
The answer we want is
\[
  EN(HH)=2+\frac14[g(H,H)+g(H,T)+g(T,H)+g(T,T)].
\]

\Knowledge{
\path{durrett_ch5_markov_property_countable};
\path{durrett_ch5_chapman_kolmogorov};
\path{durrett_ch5_expectation_operator_notation}.
}

\noindent\textbf{Solution.}
Let the target state be $HH$.  Then $g(HH)=0$.  From Exercise 5.1.3,
the transitions are
\[
  (a,b)\longrightarrow (b,H)\text{ or }(b,T),
  \qquad\text{each with probability }1/2.
\]
Write
\[
  a=g(HT),\qquad b=g(TH),\qquad c=g(TT).
\]
Using Exercise 5.2.11, the equations are
\[
  a=1+\frac12 b+\frac12 c,
\]
\[
  b=1+\frac12 g(HH)+\frac12 a
    =1+\frac12 a,
\]
and
\[
  c=1+\frac12 b+\frac12 c.
\]
The first and third equations imply $a=c$.  Hence
\[
  a=1+\frac12 b+\frac12 a,
  \qquad\text{so}\qquad a=2+b.
\]
Together with $b=1+a/2$, this gives
\[
  a=6,\qquad b=4,\qquad c=6.
\]
Thus
\[
  g(HH)=0,\qquad g(HT)=6,\qquad g(TH)=4,\qquad g(TT)=6.
\]
The first two tosses create one of the four pair states uniformly, so
\[
  EN(HH)
  =
  2+\frac14(0+6+4+6)
  =
  6.
\]

\section*{Exercise 5.2.13}

\noindent\textbf{Question.}
In this problem we consider the number of tosses $N(HHH)$ to get
$HHH$.  Let $X_n$ be the number of consecutive heads we have after $n$
tosses.  Then $X_n$ is a Markov chain with transition probability
\[
  \begin{array}{c|cccc}
    &0&1&2&3\\ \hline
  0&1/2&1/2&0&0\\
  1&1/2&0&1/2&0\\
  2&1/2&0&0&1/2\\
  3&0&0&0&1
  \end{array}
\]
Let $T_3$ be the time to reach state $3$.  Use Problem 5.2.11 to compute
$E_0T_3$.

\Knowledge{
\path{durrett_ch5_markov_property_countable};
\path{durrett_ch5_expectation_operator_notation};
\path{durrett_ch5_hitting_return_decomposition}.
}

\noindent\textbf{Solution.}
Let
\[
  g_i=E_iT_3,\qquad i=0,1,2,3.
\]
Then $g_3=0$.  By Exercise 5.2.11, for $i=0,1,2$,
\[
  g_i=1+\sum_j p(i,j)g_j.
\]
Using the displayed transition matrix,
\[
  g_0=1+\frac12g_0+\frac12g_1,
\]
\[
  g_1=1+\frac12g_0+\frac12g_2,
\]
and
\[
  g_2=1+\frac12g_0+\frac12g_3
      =1+\frac12g_0.
\]
Substitute the last equation into the equation for $g_1$:
\[
  g_1
  =
  1+\frac12g_0+\frac12\left(1+\frac12g_0\right)
  =
  \frac32+\frac34g_0.
\]
The equation for $g_0$ gives $g_0=2+g_1$, so
\[
  g_0
  =
  2+\frac32+\frac34g_0
  =
  \frac72+\frac34g_0.
\]
Therefore
\[
  \frac14g_0=\frac72,
  \qquad\text{and hence}\qquad
  E_0T_3=g_0=14.
\]

\section*{Exercise 5.3.1}

\noindent\textbf{Question.}
Suppose $y$ is recurrent and, for $k\ge0$, let
$R_k=T_y^k$ be the time of the $k$th return to $y$.  For $k\ge1$, let
$r_k=R_k-R_{k-1}$ be the $k$th interarrival time.  Use the strong Markov
property to conclude that under $P_y$, the excursion vectors
\[
  v_k=(r_k,X_{R_{k-1}},X_{R_{k-1}+1},\ldots,X_{R_k-1}),
  \qquad k\ge1,
\]
are i.i.d.

\Knowledge{
\path{durrett_ch5_strong_markov_property};
\path{durrett_ch5_hitting_return_decomposition};
\path{durrett_ch5_recurrence_visit_criterion}.
}

\noindent\textbf{Solution.}
Since $y$ is recurrent, under $P_y$ all return times $R_k$ are finite
almost surely.  The time $R_{k-1}$ is a stopping time and
$X_{R_{k-1}}=y$.  By the strong Markov property, conditional on
$\mathcal{F}_{R_{k-1}}$, the shifted process
\[
  (X_{R_{k-1}+m})_{m\ge0}
\]
has the same law as a fresh chain started from $y$, and is independent
of the past before $R_{k-1}$.

Therefore the path segment from $R_{k-1}$ up to, but not including, the
next return to $y$ has the same distribution for every $k$.  Moreover,
each such segment is conditionally independent of all earlier segments.
Thus the excursion vectors $v_1,v_2,\ldots$ are independent and
identically distributed under $P_y$.

\section*{Exercise 5.3.2}

\noindent\textbf{Question.}
Use the strong Markov property to show that
\[
  \rho_{xz}\ge \rho_{xy}\rho_{yz}.
\]

\Knowledge{
\path{durrett_ch5_strong_markov_property};
\path{durrett_ch5_hitting_return_decomposition}.
}

\noindent\textbf{Solution.}
Let $T_y=\inf\{n\ge1:X_n=y\}$ and
$T_z=\inf\{n\ge1:X_n=z\}$.  If the chain first hits $y$ and then, from
$y$, eventually hits $z$, then it has hit $z$ starting from $x$.  By the
strong Markov property at $T_y$,
\[
\begin{aligned}
  \rho_{xz}
  &=P_x(T_z<\infty)\\
  &\ge
  P_x(T_y<\infty,\ T_z\circ\theta_{T_y}<\infty)\\
  &=
  E_x\!\left[
    \1_{\{T_y<\infty\}}
    P_y(T_z<\infty)
  \right]\\
  &=\rho_{xy}\rho_{yz}.
\end{aligned}
\]

\section*{Exercise 5.3.3}

\noindent\textbf{Question.}
Show that in the Ehrenfest chain from Example 5.1.3 all states are
recurrent.

\Knowledge{
\path{durrett_ch5_ehrenfest_chain};
\path{durrett_ch5_finite_closed_recurrent_class};
\path{durrett_ch5_communication_recurrence_transfer}.
}

\noindent\textbf{Solution.}
The Ehrenfest chain has finite state space $\{0,1,\ldots,r\}$ and
transition probabilities
\[
  p(k,k+1)=\frac{r-k}{r},
  \qquad
  p(k,k-1)=\frac{k}{r}.
\]
From any state $i$ one can reach any state $j$ by moving one step at a
time upward or downward, and each required step has positive
probability.  Hence the whole finite state space is irreducible.  It is
also closed, since the chain never leaves $\{0,\ldots,r\}$.

A finite closed irreducible class consists entirely of recurrent states.
Therefore every state of the Ehrenfest chain is recurrent.

\section*{Exercise 5.3.4}

\noindent\textbf{Question.}
To probe the boundary between recurrence and transience for birth and
death chains, suppose
\[
  p_j=\frac12+\epsilon_j,\qquad
  \epsilon_j\sim Cj^{-\alpha}\quad\text{as }j\to\infty,
  \qquad
  q_j=1-p_j.
\]
Show that:
\begin{enumerate}[label=(\roman*)]
\item if $\alpha>1$, then $0$ is recurrent;
\item if $\alpha<1$, then $0$ is transient;
\item if $\alpha=1$, then $0$ is transient if $C>1/4$ and recurrent if
$C<1/4$.
\end{enumerate}

\Knowledge{
\path{durrett_ch5_birth_death_scale_function};
\path{durrett_ch5_birth_death_recurrence_test}.
}

\noindent\textbf{Solution.}
For a birth and death chain, the recurrence test is
\[
  \sum_{m=0}^\infty \prod_{j=1}^m \frac{q_j}{p_j}=\infty.
\]
Here
\[
  \frac{q_j}{p_j}
  =
  \frac{1/2-\epsilon_j}{1/2+\epsilon_j}
  =
  \frac{1-2\epsilon_j}{1+2\epsilon_j}.
\]
For small $\epsilon_j$,
\[
  \log\frac{q_j}{p_j}
  =
  -4\epsilon_j+O(\epsilon_j^2).
\]
Thus
\[
  \log\prod_{j=1}^m\frac{q_j}{p_j}
  =
  -4\sum_{j=1}^m\epsilon_j+O\!\left(\sum_{j=1}^m\epsilon_j^2\right).
\]

If $\alpha>1$, then $\sum_j \epsilon_j$ converges, and so the products
$\prod_{j=1}^m(q_j/p_j)$ converge to a positive finite limit.  The
series of products diverges, so $0$ is recurrent.

If $\alpha<1$ and $C>0$, then
\[
  \sum_{j=1}^m\epsilon_j
  \sim
  \frac{C}{1-\alpha}m^{1-\alpha}.
\]
The products are bounded above by
\[
  \exp\{-c m^{1-\alpha}\}
\]
for some $c>0$ and large $m$, so the series of products converges.  Thus
$0$ is transient.  This is the outward-drift case intended by the
statement; if $C<0$, the drift is eventually toward $0$ and the same
test gives recurrence.

Finally suppose $\alpha=1$.  Then
\[
  \sum_{j=1}^m\epsilon_j\sim C\log m,
\]
and since $\sum_j\epsilon_j^2<\infty$,
\[
  \prod_{j=1}^m\frac{q_j}{p_j}
  =
  m^{-4C+o(1)}.
\]
Therefore the product series behaves like a $p$-series with exponent
$4C$.  It diverges when $4C<1$ and converges when $4C>1$.  Hence $0$ is
recurrent if $C<1/4$ and transient if $C>1/4$.

\section*{Exercise 5.3.5}

\noindent\textbf{Question.}
Show that if we replace ``$\varphi(x)\to\infty$'' by
``$\varphi(x)\to0$'' in Theorem 5.3.8 and assume that
$\varphi(x)>0$ for $x\in F$, then we can conclude that the chain is
transient.

\Knowledge{
\path{durrett_ch5_foster_supermartingale_recurrence};
\path{durrett_ch5_recurrence_visit_criterion}.
}

\noindent\textbf{Solution.}
Assume the chain is irreducible and that
\[
  E_x\varphi(X_1)\le\varphi(x)
  \qquad x\notin F,
\]
where $F$ is finite, $\varphi\ge0$, $\varphi(x)\to0$ as $x\to\infty$,
and $\varphi>0$ on $F$.

Let
\[
  c=\min_{z\in F}\varphi(z)>0.
\]
Since $\varphi(x)\to0$, choose $x\notin F$ with $\varphi(x)<c$.  Let
\[
  \tau=\inf\{n\ge0:X_n\in F\}.
\]
Then $\varphi(X_{n\wedge\tau})$ is a nonnegative supermartingale.  If
the chain were recurrent, irreducibility would imply $P_x(\tau<\infty)=1$.
By optional stopping and Fatou's lemma,
\[
  \varphi(x)
  \ge
  E_x\varphi(X_\tau)
  \ge
  cP_x(\tau<\infty)
  =
  c,
\]
contradicting $\varphi(x)<c$.

Thus the irreducible chain cannot be recurrent.  In an irreducible
chain, all states are either recurrent or transient, so the chain is
transient.

\section*{Exercise 5.3.6}

\noindent\textbf{Question.}
Let $X_n$ be a birth and death chain with
\[
  p_j-\frac12\sim \frac{C}{j}
  \quad\text{as }j\to\infty,
  \qquad q_j=1-p_j.
\]
\begin{enumerate}[label=(\roman*)]
\item Show that if $C<1/4$, then we can pick $\alpha>0$ so that
$\varphi(x)=x^\alpha$ satisfies the hypotheses of Theorem 5.3.8.
\item Show that when $C>1/4$, we can take $\alpha<0$ and apply Exercise
5.3.5.
\end{enumerate}

\Knowledge{
\path{durrett_ch5_foster_supermartingale_recurrence};
\path{durrett_ch5_birth_death_recurrence_test}.
}

\noindent\textbf{Solution.}
For large $j$, write $p_j=1/2+C/j+o(1/j)$ and
$q_j=1/2-C/j+o(1/j)$.  For $\varphi(x)=x^\alpha$, Taylor expansion gives
\[
  (j+1)^\alpha-j^\alpha
  =
  \alpha j^{\alpha-1}
  +\frac{\alpha(\alpha-1)}{2}j^{\alpha-2}
  +o(j^{\alpha-2}),
\]
and
\[
  (j-1)^\alpha-j^\alpha
  =
  -\alpha j^{\alpha-1}
  +\frac{\alpha(\alpha-1)}{2}j^{\alpha-2}
  +o(j^{\alpha-2}).
\]
Hence
\[
\begin{aligned}
  E_j\varphi(X_1)-\varphi(j)
  &=
  p_j\bigl((j+1)^\alpha-j^\alpha\bigr)
  +q_j\bigl((j-1)^\alpha-j^\alpha\bigr)\\
  &=
  \alpha\left(2C+\frac{\alpha-1}{2}\right)j^{\alpha-2}
  +o(j^{\alpha-2}).
\end{aligned}
\]

If $C<1/4$, choose $\alpha>0$ small enough that
\[
  2C+\frac{\alpha-1}{2}<0,
  \qquad\text{equivalently}\qquad
  \alpha<1-4C.
\]
Then $E_j\varphi(X_1)\le\varphi(j)$ for all sufficiently large $j$, and
$\varphi(j)=j^\alpha\to\infty$.  Taking the exceptional small states as
the finite set $F$, Theorem 5.3.8 implies recurrence.

If $C>1/4$, then $1-4C<0$.  Choose
\[
  1-4C<\alpha<0.
\]
Now $\alpha<0$ and
$2C+(\alpha-1)/2>0$, so their product is negative.  Therefore
$E_j\varphi(X_1)\le\varphi(j)$ for all sufficiently large $j$, while
$\varphi(j)=j^\alpha\to0$ and $\varphi>0$ on every finite set.  Exercise
5.3.5 applies and gives transience.

\section*{Exercise 5.3.7}

\noindent\textbf{Question.}
A function $f$ is said to be superharmonic if
\[
  f(x)\ge \sum_y p(x,y)f(y),
\]
or equivalently $f(X_n)$ is a supermartingale.  Suppose $p$ is
irreducible.  Show that $p$ is recurrent if and only if every
nonnegative superharmonic function is constant.

\Knowledge{
\path{durrett_ch5_recurrence_visit_criterion};
\path{durrett_ch5_communication_recurrence_transfer};
\path{durrett_ch5_foster_supermartingale_recurrence}.
}

\noindent\textbf{Solution.}
First suppose the chain is recurrent and let $f\ge0$ be superharmonic.
For states $x,y$, irreducibility and recurrence imply
$P_x(T_y<\infty)=1$.  Since $f(X_n)$ is a nonnegative supermartingale,
optional stopping at $n\wedge T_y$ and Fatou's lemma give
\[
  f(x)
  \ge
  E_x f(X_{T_y})
  =
  f(y).
\]
Interchanging $x$ and $y$ gives $f(y)\ge f(x)$, so $f(x)=f(y)$.  Since
$x,y$ were arbitrary, $f$ is constant.

Conversely, suppose the irreducible chain is transient.  Fix a state
$a$ and define
\[
  f(x)=P_x(T_a<\infty).
\]
Then $f\ge0$ and $f(a)=1$.  For $x\ne a$, first-step decomposition gives
\[
  f(x)=\sum_y p(x,y)f(y).
\]
At $x=a$,
\[
  \sum_y p(a,y)f(y)=P_a(T_a<\infty)=\rho_{aa}<1=f(a),
\]
because $a$ is transient.  Thus $f$ is nonnegative and superharmonic.
It is not constant, since $f(a)=1$ but the displayed strict inequality
would be impossible for a constant function.  Therefore if every
nonnegative superharmonic function is constant, the chain must be
recurrent.

\section*{Exercise 5.3.8}

\noindent\textbf{Question.}
$M/M/\infty$ queue.  Consider a telephone system with an infinite number
of lines.  Let $X_n$ be the number of lines in use at time $n$, and
suppose
\[
  X_{n+1}=\sum_{m=1}^{X_n}\xi_{n,m}+Y_{n+1},
\]
where the $\xi_{n,m}$ are i.i.d. with
$P(\xi_{n,m}=1)=p$ and $P(\xi_{n,m}=0)=1-p$, and $Y_n$ is an independent
i.i.d. sequence of Poisson mean $\lambda$ random variables.  In words,
for each conversation we flip a coin with probability $p$ of heads to
see if it continues for another minute.  Meanwhile, a Poisson mean
$\lambda$ number of conversations start between time $n$ and $n+1$.  Use
Theorem 5.3.8 with $\varphi(x)=x$ to show that the chain is recurrent
for any $p<1$.

\Knowledge{
\path{durrett_ch5_foster_supermartingale_recurrence};
\path{durrett_ch5_markov_property_countable}.
}

\noindent\textbf{Solution.}
Given $X_n=x$, the number of continuing old conversations has
binomial$(x,p)$ distribution, and the number of new conversations has
Poisson mean $\lambda$.  Therefore
\[
  E_xX_1=px+\lambda.
\]
Let $\varphi(x)=x$.  Since $p<1$, choose a finite set
\[
  F=\{0,1,\ldots,K\}
\]
with $K$ large enough that
\[
  px+\lambda\le x
  \qquad\text{for all }x>K.
\]
Equivalently, take $K\ge \lambda/(1-p)$.  Then for $x\notin F$,
\[
  E_x\varphi(X_1)=E_xX_1=px+\lambda\le x=\varphi(x).
\]
Also $\varphi(x)=x\to\infty$.

If $\lambda>0$, the chain is irreducible on $\{0,1,2,\ldots\}$: from
any $x$ all current conversations can end with positive probability, and
the Poisson arrivals can produce any desired number of new
conversations with positive probability.  Theorem 5.3.8 therefore
implies the chain is recurrent.  If $\lambda=0$, state $0$ is absorbing
and every initial state reaches it with probability one when $p<1$, so
the recurrent class is again evident.

\section*{Exercise 5.4.1}

\noindent\textbf{Question.}
For a random walk on $\mathbb{R}$, show that there are only four
possibilities, one of which has probability one:
\begin{enumerate}[label=(\roman*)]
\item $S_n=0$ for all $n$.
\item $S_n\to\infty$.
\item $S_n\to-\infty$.
\item $-\infty=\liminf S_n<\limsup S_n=\infty$.
\end{enumerate}

\Knowledge{
\path{durrett_ch5_random_walk_chain};
\path{durrett_ch5_random_walk_recurrent_values};
\path{durrett_ch5_random_walk_recurrence_equivalences}.
}

\noindent\textbf{Solution.}
Let
\[
  M^+=\sup_{n\ge0}S_n,
  \qquad
  M^-=\sup_{n\ge0}(-S_n)=-\inf_{n\ge0}S_n.
\]
The events $\{M^+=\infty\}$ and $\{M^-=\infty\}$ are tail events for
the iid increments, so each has probability $0$ or $1$.

If both have probability $1$, then
\[
  \limsup_{n\to\infty}S_n=\infty,
  \qquad
  \liminf_{n\to\infty}S_n=-\infty,
\]
which is case (iv).

If $\Prob(M^+=\infty)=1$ and $\Prob(M^-=\infty)=0$, then the walk is
unbounded above but bounded below.  The standard ladder-time argument
then implies $S_n\to\infty$: if the walk returned below a fixed level
infinitely often while also being unbounded above, then independent
future excursions from those return times would, with probability one,
eventually create arbitrarily low new records, contradicting
$M^-<\infty$.  This is case (ii).  By applying the same argument to
$-S_n$, the case $\Prob(M^+=\infty)=0$ and $\Prob(M^-=\infty)=1$ gives
$S_n\to-\infty$, which is case (iii).

It remains to consider the case in which both $M^+$ and $M^-$ are finite
a.s.  Then the random walk is a.s. bounded.  If the increment law were
not concentrated at $0$, then for some $\varepsilon>0$ either
$\Prob(\xi_1>\varepsilon)>0$ or $\Prob(\xi_1<-\varepsilon)>0$.  In the
first case, arbitrarily long runs of increments larger than
$\varepsilon$ occur a.s.; in the second case, arbitrarily long runs of
increments smaller than $-\varepsilon$ occur a.s.  Either kind of run
contradicts boundedness of the partial sums.  Hence $\xi_1=0$ a.s., and
$S_n=0$ for all $n$, which is case (i).

\section*{Exercise 5.4.2}

\noindent\textbf{Question.}
Ladder variables.  Let
\[
  \alpha(\omega)=\inf\{n:\omega_1+\cdots+\omega_n>0\},
\]
where $\inf\varnothing=\infty$, and set $\alpha(\Delta)=\infty$.  Let
$\alpha_0=0$ and
\[
  \alpha_k(\omega)=\alpha_{k-1}(\omega)+\alpha(\theta_{\alpha_{k-1}}\omega),
  \qquad k\ge1.
\]
At time $\alpha_k$, the random walk is at a record high value.
\begin{enumerate}[label=(\roman*)]
\item If $P(\alpha<\infty)<1$, then $P(\sup S_n<\infty)=1$.
\item If $P(\alpha<\infty)=1$, then $P(\sup S_n=\infty)=1$.
\end{enumerate}

\Knowledge{
\path{durrett_ch5_strong_markov_property};
\path{durrett_ch5_random_walk_chain};
\path{durrett_ch5_hitting_return_decomposition}.
}

\noindent\textbf{Solution.}
Let $p=P(\alpha<\infty)$.  By the strong Markov property applied at the
successive ladder times,
\[
  P(\alpha_k<\infty)=p^k.
\]
Indeed, each time a new strict record is reached, the future increments
are independent of the past and have the same law as the original
increments.

If $p<1$, then $P(\alpha_k<\infty)=p^k\to0$.  Hence with probability
one, only finitely many strict ascending ladder times occur.  After the
last one, the walk never rises above the last record value, so
\[
  \sup_{n\ge0}S_n<\infty
  \qquad\text{a.s.}
\]

If $p=1$, then every $\alpha_k$ is finite a.s.  The ladder height
increments
\[
  S_{\alpha_k}-S_{\alpha_{k-1}},
  \qquad k\ge1,
\]
are iid and strictly positive.  Their partial sums therefore diverge to
$\infty$ a.s.  Since $S_{\alpha_k}$ is the $k$th strict record value,
\[
  \sup_{n\ge0}S_n=\infty
  \qquad\text{a.s.}
\]

\section*{Exercise 5.4.3}

\noindent\textbf{Question.}
Suppose $f$ is superharmonic on $\mathbb{R}^d$, i.e.
\[
  f(x)\ge
  \frac{1}{|B(x,r)|}\int_{B(x,r)} f(y)\,dy.
\]
Let $\xi_1,\xi_2,\ldots$ be i.i.d. uniform on $B(0,1)$, and define
$S_n$ by $S_n=S_{n-1}+\xi_n$ for $n\ge1$ and $S_0=x$.
\begin{enumerate}[label=(\roman*)]
\item Show that $X_n=f(S_n)$ is a supermartingale.
\item Use this to conclude that in $d\le2$, nonnegative superharmonic
functions must be constant.  The example $f(x)=|x|^{2-d}$ shows this is
false in $d>2$.
\end{enumerate}

\Knowledge{
\path{durrett_ch5_chung_fuchs_one_dimension};
\path{durrett_ch5_two_dimensional_normal_recurrence};
\path{durrett_ch5_random_walk_recurrent_values};
\path{durrett_ch5_foster_supermartingale_recurrence}.
}

\noindent\textbf{Solution.}
Since $\xi_{n+1}$ is uniform on $B(0,1)$ and independent of
$\mathcal{F}_n$,
\[
\begin{aligned}
  E(f(S_{n+1})\mid\mathcal{F}_n)
  &=
  \frac{1}{|B(0,1)|}\int_{B(0,1)} f(S_n+z)\,dz\\
  &=
  \frac{1}{|B(S_n,1)|}\int_{B(S_n,1)} f(y)\,dy\\
  &\le f(S_n).
\end{aligned}
\]
Thus $f(S_n)$ is a supermartingale whenever the expectations are
well-defined; for nonnegative $f$, this is understood in the usual
extended nonnegative sense.

Now assume $d\le2$ and $f\ge0$.  The increment distribution is centered
and nondegenerate.  In $d=1$, Chung-Fuchs recurrence applies; in $d=2$,
the central limit theorem gives a nondegenerate normal limit after
$\sqrt n$ scaling, so the planar recurrence theorem applies.  Hence the
random walk is recurrent, and because the increment law has positive
density on a ball, every nonempty open ball is visited infinitely often
from every starting point.

The nonnegative supermartingale $f(S_n)$ converges a.s.  If two points
could have different superharmonic values, recurrence would force the
walk to visit neighborhoods of both points infinitely often, contradicting
the existence of a single almost sure limit.  Equivalently, this is the
general-state analogue of the recurrent-chain fact that every
nonnegative superharmonic function is constant.  Therefore every
nonnegative superharmonic function on $\mathbb{R}^d$ is constant for
$d\le2$.

\section*{Exercise 5.4.4}

\noindent\textbf{Question.}
Suppose $h$ is harmonic on $\mathbb{R}^d$, i.e.
\[
  h(x)=
  \frac{1}{|B(x,r)|}\int_{B(x,r)} h(y)\,dy.
\]
Let $\xi_1,\xi_2,\ldots$ be i.i.d. uniform on $B(0,1)$, and define
$S_n$ by $S_n=S_{n-1}+\xi_n$ for $n\ge1$ and $S_0=x$.
\begin{enumerate}[label=(\roman*)]
\item Show that $X_n=h(S_n)$ is a martingale.
\item Use this to conclude that in any dimension, bounded harmonic
functions must be constant.  The Hewitt-Savage $0$-$1$ law will be
useful here.
\end{enumerate}

\Knowledge{
\path{durrett_ch5_random_walk_chain};
\path{durrett_ch5_markov_property_countable};
\path{durrett_ch5_random_walk_recurrence_equivalences}.
}

\noindent\textbf{Solution.}
As in the previous exercise, conditioning on $\mathcal{F}_n$ gives
\[
\begin{aligned}
  E(h(S_{n+1})\mid\mathcal{F}_n)
  &=
  \frac{1}{|B(0,1)|}\int_{B(0,1)} h(S_n+z)\,dz\\
  &=
  \frac{1}{|B(S_n,1)|}\int_{B(S_n,1)} h(y)\,dy\\
  &=h(S_n).
\end{aligned}
\]
Thus $h(S_n)$ is a martingale.

Assume now that $h$ is bounded.  Then $h(S_n)$ is a bounded martingale,
so it converges a.s. to a random variable $L_x$ when $S_0=x$.  If we
permute finitely many increments, then all sufficiently late partial
sums are unchanged, since the total sum of those finitely many
increments is unchanged.  Hence $L_x$ is invariant under finite
permutations of the iid sequence $(\xi_i)$.  By the Hewitt-Savage
$0$-$1$ law, $L_x$ is a.s. constant.  Since
\[
  h(x)=E_x h(S_n)=E_x L_x,
\]
this constant is $h(x)$.

It remains to see that the constant does not depend on $x$.  Fix
$k\ge1$.  By the Markov property,
\[
  L_x=L_{S_k}\quad\text{a.s.}
\]
The distribution of $S_k$ has a density that is positive on the ball
$B(x,k)$, so $h(S_k)=h(x)$ a.s. implies that $h(z)=h(x)$ for Lebesgue
almost every $z\in B(x,k)$.  Since $k$ is arbitrary, $h$ is a.e. equal
to the constant $h(x)$ on all of $\mathbb{R}^d$.  Applying the harmonic
mean-value identity once more yields equality at every point.  Therefore
every bounded harmonic function on $\mathbb{R}^d$ is constant.

\section*{Exercise 5.5.1}

\noindent\textbf{Question.}
Find the stationary distribution for the Bernoulli-Laplace model of
diffusion from Exercise 5.1.5.

\Knowledge{
\path{durrett_ch5_stationary_measure_definition};
\path{durrett_ch5_reversibility_detailed_balance};
\path{durrett_ch5_birth_death_stationary_distribution}.
}

\noindent\textbf{Solution.}
The state $i$ is the number of black balls in the left urn.  Since
$b\le m$, the state space is $\{0,1,\ldots,b\}$.  From Exercise 5.1.5,
\[
  p(i,i+1)=\frac{(m-i)(b-i)}{m^2},
  \qquad
  p(i,i-1)=\frac{i(m-b+i)}{m^2}.
\]
At equilibrium, the left urn should look like a uniformly chosen set of
$m$ balls from the $2m$ total balls.  Hence the number of black balls in
the left urn has the hypergeometric distribution
\[
  \pi(i)=
  \frac{\binom{b}{i}\binom{2m-b}{m-i}}{\binom{2m}{m}},
  \qquad 0\le i\le b.
\]
We can verify this by detailed balance.  For $0\le i<b$,
\[
  \frac{\pi(i+1)}{\pi(i)}
  =
  \frac{b-i}{i+1}\cdot
  \frac{m-i}{m-b+i+1}.
\]
Also
\[
  \frac{p(i,i+1)}{p(i+1,i)}
  =
  \frac{(m-i)(b-i)}{(i+1)(m-b+i+1)}.
\]
Thus
\[
  \pi(i)p(i,i+1)=\pi(i+1)p(i+1,i),
\]
so $\pi$ satisfies detailed balance and is stationary.

\section*{Exercise 5.5.2}

\noindent\textbf{Question.}
Let
\[
  w_{xy}=P_x(T_y<T_x).
\]
Show that
\[
  \mu_x(y)=\frac{w_{xy}}{w_{yx}}.
\]

\Knowledge{
\path{durrett_ch5_stationary_measure_from_returns};
\path{durrett_ch5_hitting_return_decomposition};
\path{durrett_ch5_strong_markov_property}.
}

\noindent\textbf{Solution.}
Assume $x$ is recurrent and $y$ is in the same recurrent class.  The
quantity $\mu_x(y)$ is the expected number of visits to $y$ before the
first positive return to $x$, starting from $x$.

First, the chain must hit $y$ before returning to $x$, an event with
probability $w_{xy}$.  Once it reaches $y$, count visits to $y$ before
hitting $x$.  Starting from $y$, each visit to $y$ starts a new trial:
with probability
\[
  w_{yx}=P_y(T_x<T_y),
\]
the chain hits $x$ before returning again to $y$; otherwise it returns
to $y$ first and the experiment restarts.  Therefore the number of
visits to $y$ before hitting $x$, counting the initial visit to $y$, has
geometric mean $1/w_{yx}$.

By the strong Markov property at the first hit of $y$,
\[
  \mu_x(y)=w_{xy}\cdot \frac1{w_{yx}}
  =
  \frac{w_{xy}}{w_{yx}}.
\]

\section*{Exercise 5.5.3}

\noindent\textbf{Question.}
Show that if $p$ is irreducible and recurrent, then
\[
  \mu_x(y)\mu_y(z)=\mu_x(z).
\]

\Knowledge{
\path{durrett_ch5_stationary_measure_from_returns};
\path{durrett_ch5_stationary_measure_uniqueness_recurrent}.
}

\noindent\textbf{Solution.}
For each recurrent state $a$, the cycle trick defines a stationary
measure $\mu_a$ with normalization $\mu_a(a)=1$.  Since the chain is
irreducible and recurrent, the stationary measure is unique up to a
constant multiple.

Thus $\mu_x$ and $\mu_y$ are multiples of one another.  Since
$\mu_y(y)=1$, evaluating the proportionality at $y$ gives
\[
  \mu_x=\mu_x(y)\mu_y.
\]
Evaluating this identity at $z$ yields
\[
  \mu_x(z)=\mu_x(y)\mu_y(z),
\]
which is the desired formula.

\section*{Exercise 5.5.4}

\noindent\textbf{Question.}
Suppose $p$ is irreducible and positive recurrent.  Then
\[
  E_xT_y<\infty
  \qquad\text{for all }x,y.
\]

\Knowledge{
\path{durrett_ch5_positive_recurrence_equivalences};
\path{durrett_ch5_positive_recurrence_return_time};
\path{durrett_ch5_strong_markov_property}.
}

\noindent\textbf{Solution.}
Since the chain is irreducible and positive recurrent, every state is
positive recurrent.  In particular,
\[
  E_yT_y<\infty.
\]
We show that this forces $E_xT_y<\infty$ for all $x$.

By irreducibility, choose a path from $y$ to $x$ with positive
probability.  Taking the first such path with minimal length, it avoids
$y$ until its endpoint $x$.  Hence there are $K<\infty$ and $a>0$ such
that
\[
  P_y(X_K=x,\ T_y>K)\ge a.
\]
If $E_xT_y=\infty$, then on the event above the return time to $y$
would have infinite conditional mean after time $K$.  By the strong
Markov property,
\[
  E_yT_y
  \ge
  a\,(K+E_xT_y)
  =
  \infty,
\]
contradicting positive recurrence of $y$.  Therefore
$E_xT_y<\infty$.

\section*{Exercise 5.5.5}

\noindent\textbf{Question.}
Suppose $p$ is irreducible and has a stationary measure $\mu$ with
\[
  \sum_x \mu(x)=\infty.
\]
Then $p$ is not positive recurrent.

\Knowledge{
\path{durrett_ch5_stationary_measure_uniqueness_recurrent};
\path{durrett_ch5_positive_recurrence_equivalences};
\path{durrett_ch5_positive_recurrence_return_time}.
}

\noindent\textbf{Solution.}
Suppose, for contradiction, that $p$ is positive recurrent.  Since the
chain is irreducible, positive recurrence is equivalent to the existence
of a stationary distribution $\pi$.

For an irreducible positive recurrent chain, every stationary measure is
a constant multiple of $\pi$.  Hence there is a constant $c\ge0$ such
that
\[
  \mu(x)=c\pi(x)
  \qquad\text{for all }x.
\]
It follows that
\[
  \sum_x\mu(x)=c\sum_x\pi(x)=c<\infty,
\]
contradicting the hypothesis that $\mu$ has infinite total mass.
Therefore $p$ is not positive recurrent.

\section*{Exercise 5.5.6}

\noindent\textbf{Question.}
Use Theorems 5.5.7 and 5.5.9 to show that for simple random walk:
\begin{enumerate}[label=(\roman*)]
\item the expected number of visits to $k$ between successive visits to
$0$ is $1$ for all $k$;
\item if we start from $k$, the expected number of visits to $k$ before
hitting $0$ is $2k$.
\end{enumerate}

\Knowledge{
\path{durrett_ch5_stationary_measure_from_returns};
\path{durrett_ch5_stationary_measure_uniqueness_recurrent};
\path{durrett_ch5_random_walk_chain}.
}

\noindent\textbf{Solution.}
For simple symmetric random walk on $\mathbb{Z}$, the counting measure
\[
  \nu(j)\equiv1
\]
is stationary, because the transition matrix is doubly stochastic.  The
chain is irreducible and recurrent.  By uniqueness of stationary
measures up to constants, the cycle measure $\mu_0$ from Theorem 5.5.7
must be a constant multiple of counting measure.  Since $\mu_0(0)=1$,
that constant is $1$.  Therefore
\[
  \mu_0(k)=1
  \qquad\text{for every }k\in\mathbb{Z}.
\]
This says exactly that the expected number of visits to $k$ between
successive visits to $0$ is $1$.

Now start from $k>0$ and count visits to $k$ before hitting $0$,
including the visit at time $0$.  Each time the walk is at $k$, the
trial ends successfully if the walk hits $0$ before returning to $k$.
The probability of this success is
\[
  P_k(T_0<T_k).
\]
To hit $0$ before returning to $k$, the first step must be to $k-1$,
which has probability $1/2$, and then the gambler's ruin probability of
hitting $0$ before $k$ starting from $k-1$ is $1/k$.  Hence
\[
  P_k(T_0<T_k)=\frac{1}{2k}.
\]
The number of visits to $k$ before hitting $0$ is geometric with this
success probability, so its mean is
\[
  \frac{1}{P_k(T_0<T_k)}=2k.
\]
By symmetry the analogous value for $k<0$ is $2|k|$.

\section*{Exercise 5.5.7}

\noindent\textbf{Question.}
Let $\xi_1,\xi_2,\ldots$ be i.i.d. with
$P(\xi_m=k)=a_{k+1}$ for $k\ge -1$, let
\[
  S_n=x+\xi_1+\cdots+\xi_n,
  \qquad x\ge0,
\]
and let
\[
  X_n=S_n+\left(\min_{m\le n}S_m\right)^-.
\]
Equation (5.1.1) shows that $X_n$ has the same distribution as the
$M/G/1$ queue starting from $X_0=x$.  Use this representation to
conclude that if $\mu=\sum_k k a_k<1$, then as $n\to\infty$,
\[
  \frac1n
  \left|\{m\le n:X_{m-1}=0,\ \xi_m=-1\}\right|
  \longrightarrow
  1-\mu
  \qquad\text{a.s.}
\]
and hence $\pi(0)=(1-\mu)/a_0$ as proved earlier.

\Knowledge{
\path{durrett_ch5_mg1_queue_chain};
\path{durrett_ch5_mg1_stability_stationarity};
\path{durrett_ch5_stationary_measure_definition}.
}

\noindent\textbf{Solution.}
The recursion for the queue is
\[
  X_m=(X_{m-1}+\xi_m)^+.
\]
Since $\xi_m\ge -1$, reflection occurs exactly on the event
\[
  \{X_{m-1}=0,\ \xi_m=-1\}.
\]
Let
\[
  L_n=\left|\{m\le n:X_{m-1}=0,\ \xi_m=-1\}\right|.
\]
Then the reflected-walk identity is
\[
  X_n=x+\sum_{m=1}^n\xi_m+L_n.
\]
Since $E\xi_1=\mu-1<0$, the strong law of large numbers gives
\[
  \frac1n\sum_{m=1}^n\xi_m\longrightarrow \mu-1.
\]
The explicit reflection representation also gives $X_n/n\to0$ a.s.:
the walk $S_n$ has negative linear drift, and subtracting its running
minimum leaves only a sublinear reflected height.  Dividing the identity
above by $n$ yields
\[
  \frac{L_n}{n}
  =
  \frac{X_n}{n}
  -\frac{x}{n}
  -\frac1n\sum_{m=1}^n\xi_m
  \longrightarrow
  1-\mu
  \qquad\text{a.s.}
\]

In stationarity, the event
$\{X_{m-1}=0,\xi_m=-1\}$ has probability $\pi(0)a_0$, because $\xi_m$ is
independent of the past and $P(\xi_m=-1)=a_0$.  The long-run frequency
therefore equals $\pi(0)a_0$.  Comparing with the a.s. limit above gives
\[
  \pi(0)a_0=1-\mu,
  \qquad\text{so}\qquad
  \pi(0)=\frac{1-\mu}{a_0}.
\]

\section*{Exercise 5.5.8}

\noindent\textbf{Question.}
$M/M/\infty$ queue.  In this chain, introduced in Exercise 5.3.8,
\[
  X_{n+1}=\sum_{m=1}^{X_n}\xi_{n,m}+Y_{n+1},
\]
where the $\xi_{n,m}$ are i.i.d. Bernoulli with mean $p$ and
$Y_{n+1}$ is an independent Poisson mean $\lambda$ random variable.
Compute the distribution at time $1$ when $X_0$ is Poisson$(\theta)$
and use the result to find the stationary distribution.

\Knowledge{
\path{durrett_ch5_stationary_measure_definition};
\path{durrett_ch5_foster_supermartingale_recurrence}.
}

\noindent\textbf{Solution.}
If $X_0\sim\operatorname{Poisson}(\theta)$, then thinning a Poisson
random variable shows that the number of old conversations that survive
one more minute is
\[
  \sum_{m=1}^{X_0}\xi_{0,m}
  \sim
  \operatorname{Poisson}(p\theta).
\]
This random variable is independent of the new arrivals
$Y_1\sim\operatorname{Poisson}(\lambda)$.  Therefore
\[
  X_1\sim\operatorname{Poisson}(p\theta+\lambda).
\]

For stationarity we need $p\theta+\lambda=\theta$, so
\[
  \theta=\frac{\lambda}{1-p}.
\]
Thus the stationary distribution is
\[
  \pi(j)=
  e^{-\lambda/(1-p)}
  \frac{(\lambda/(1-p))^j}{j!},
  \qquad j=0,1,2,\ldots.
\]

\section*{Exercise 5.5.9}

\noindent\textbf{Question.}
Let $X_n\ge0$ be a Markov chain and suppose
\[
  E_xX_1\le x-\epsilon
  \qquad\text{for }x>K,
\]
where $\epsilon>0$.  Let
\[
  Y_n=X_n+n\epsilon,
  \qquad
  \tau=\inf\{n:X_n\le K\}.
\]
Then $Y_{n\wedge\tau}$ is a positive supermartingale and the optional
stopping theorem implies $E_x\tau\le x/\epsilon$.  With assumptions
about the behavior of the chain starting from points $x\le K$, this
leads to conditions for positive recurrence.

\Knowledge{
\path{durrett_ch5_foster_supermartingale_recurrence};
\path{durrett_ch5_positive_recurrence_equivalences};
\path{durrett_ch5_hitting_return_decomposition}.
}

\noindent\textbf{Solution.}
For $n<\tau$, we have $X_n>K$, so
\[
  E(X_{n+1}\mid\mathcal{F}_n)\le X_n-\epsilon.
\]
Hence
\[
  E(Y_{n+1}\mid\mathcal{F}_n)
  =
  E(X_{n+1}\mid\mathcal{F}_n)+(n+1)\epsilon
  \le
  X_n+n\epsilon
  =
  Y_n.
\]
After $\tau$ the stopped process is constant, so
$Y_{n\wedge\tau}$ is a nonnegative supermartingale.

By optional stopping,
\[
  E_xY_{n\wedge\tau}\le Y_0=x.
\]
Since $X_{n\wedge\tau}\ge0$,
\[
  \epsilon E_x(n\wedge\tau)\le x.
\]
Letting $n\to\infty$ and using monotone convergence gives
\[
  E_x\tau\le \frac{x}{\epsilon}.
\]

Let $C=\{x:x\le K\}$.  If $C$ is finite and the chain is irreducible on
a countable state space, this drift estimate is the key Foster
condition for positive recurrence.  Indeed, from every outside state the
expected time to hit $C$ is finite.  Starting from the finite set $C$,
irreducibility gives a fixed state $a\in C$ that can be reached with
positive probability in finitely many steps; failed attempts return to
$C$ in finite mean time by the drift estimate.  A geometric-trials
argument then gives $E_aT_a<\infty$.  Thus $a$ is positive recurrent,
and by irreducibility all states are positive recurrent.

\section*{Exercise 5.6.1}

\noindent\textbf{Question.}
Suppose $S=\{0,1\}$ and
\[
  p=
  \begin{pmatrix}
    1-\alpha & \alpha\\
    \beta & 1-\beta
  \end{pmatrix}.
\]
Use induction to show that
\[
  P_\mu(X_n=0)
  =
  \frac{\beta}{\alpha+\beta}
  +(1-\alpha-\beta)^n
  \left\{\mu(0)-\frac{\beta}{\alpha+\beta}\right\}.
\]

\Knowledge{
\path{durrett_ch5_stationary_measure_definition};
\path{durrett_ch5_aperiodic_convergence_theorem}.
}

\noindent\textbf{Solution.}
Let
\[
  q_n=P_\mu(X_n=0).
\]
Then
\[
  q_{n+1}
  =
  q_n(1-\alpha)+(1-q_n)\beta
  =
  \beta+(1-\alpha-\beta)q_n.
\]
The fixed point of this affine recursion is
\[
  q_*=\frac{\beta}{\alpha+\beta},
\]
since $q_*=\beta+(1-\alpha-\beta)q_*$.  Therefore
\[
  q_{n+1}-q_*=(1-\alpha-\beta)(q_n-q_*).
\]
Iterating from $q_0=\mu(0)$ gives
\[
  q_n-q_*=(1-\alpha-\beta)^n(\mu(0)-q_*),
\]
which is the desired formula.

\section*{Exercise 5.6.2}

\noindent\textbf{Question.}
Show that if $S$ is finite and $p$ is irreducible and aperiodic, then
there is an $m$ such that
\[
  p^m(x,y)>0
  \qquad\text{for all }x,y.
\]

\Knowledge{
\path{durrett_ch5_period_definition};
\path{durrett_ch5_aperiodic_convergence_theorem};
\path{durrett_ch5_finite_closed_recurrent_class}.
}

\noindent\textbf{Solution.}
Because the chain is irreducible, for each pair $x,y$ there is a
$k(x,y)$ such that
\[
  p^{k(x,y)}(x,y)>0.
\]
Because the chain is aperiodic, Lemma 5.6.5 from the textbook says that
for each $y$ there is an $M(y)$ such that
\[
  p^\ell(y,y)>0
  \qquad\text{for all }\ell\ge M(y).
\]
The state space is finite, so choose
\[
  m\ge \max_{x,y}\{k(x,y)+M(y)\}.
\]
Then for every $x,y$,
\[
  p^m(x,y)
  \ge
  p^{k(x,y)}(x,y)\,p^{m-k(x,y)}(y,y)
  >0.
\]

\section*{Exercise 5.6.3}

\noindent\textbf{Question.}
Show that if $S$ is finite, $p$ is irreducible and aperiodic, and $T$
is the coupling time defined in the proof of Theorem 5.6.6, then
\[
  P(T>n)\le Cr^n
\]
for some $r<1$ and $C<\infty$.  Thus convergence to equilibrium occurs
exponentially rapidly in this case.  Hint: first consider the case in
which $p(x,y)>0$ for all $x,y$ and reduce the general case to this one
by looking at a power of $p$.

\Knowledge{
\path{durrett_ch5_aperiodic_convergence_theorem};
\path{durrett_ch5_coupling_for_finite_chains};
\path{durrett_ch5_period_definition}.
}

\noindent\textbf{Solution.}
First suppose $p(x,y)>0$ for all $x,y$.  Since $S$ is finite, choose a
state $a$ and set
\[
  \eta=\min_{x\in S}p(x,a)>0.
\]
For the independent-coordinate coupling used in Theorem 5.6.6, from any
current pair $(x,y)$ the probability that both coordinates jump to $a$
in the next step is at least $\eta^2$.  Hence
\[
  P(T>n+1\mid T>n)\le 1-\eta^2,
\]
and therefore
\[
  P(T>n)\le(1-\eta^2)^n.
\]

In the general finite irreducible aperiodic case, Exercise 5.6.2 gives
an integer $m$ such that
\[
  p^m(x,y)>0
  \qquad\text{for all }x,y.
\]
Apply the previous argument to the $m$-step skeleton chain.  There is
$\eta>0$ such that, at each block of $m$ steps, the coupled coordinates
meet with probability at least $\eta^2$.  Thus
\[
  P(T>km)\le(1-\eta^2)^k.
\]
For arbitrary $n$, take $k=\lfloor n/m\rfloor$ and write this as
\[
  P(T>n)\le C r^n
\]
for suitable constants $C<\infty$ and $r=(1-\eta^2)^{1/m}<1$.

\section*{Exercise 5.6.4}

\noindent\textbf{Question.}
For any transition matrix $p$, define
\[
  \alpha_n
  =
  \sup_{i,j}
  \frac12\sum_k |p^n(i,k)-p^n(j,k)|.
\]
The factor $1/2$ is there because for any $i,j$ we can define random
variables $X$ and $Y$ such that
$P(X=k)=p^n(i,k)$, $P(Y=k)=p^n(j,k)$, and
\[
  P(X\ne Y)
  =
  \frac12\sum_k |p^n(i,k)-p^n(j,k)|.
\]
Show that
\[
  \alpha_{m+n}\le \alpha_n\alpha_m.
\]

\Knowledge{
\path{durrett_ch5_coupling_for_finite_chains};
\path{durrett_ch5_aperiodic_convergence_theorem}.
}

\noindent\textbf{Solution.}
We use the coupling interpretation.  Start two copies of the chain from
$i$ and $j$.  Couple their positions after $m$ steps as efficiently as
possible, so that
\[
  P(X_m\ne Y_m)
  \le
  \alpha_m.
\]
If $X_m=Y_m$, then use the same future randomness for the two chains,
so they remain together.  If $X_m\ne Y_m$, couple the next $n$ steps as
efficiently as possible; conditional on the two different states at time
$m$, the chance they fail to couple after another $n$ steps is at most
$\alpha_n$.

Therefore
\[
  P(X_{m+n}\ne Y_{m+n})\le \alpha_m\alpha_n.
\]
Taking the best possible coupling interpretation of total variation and
then the supremum over $i,j$ gives
\[
  \alpha_{m+n}\le\alpha_m\alpha_n.
\]

\section*{Exercise 5.6.5}

\noindent\textbf{Question.}
Strong law for additive functionals.  Suppose $p$ is irreducible and has
stationary distribution $\pi$.  Let $f$ satisfy
\[
  \sum_y |f(y)|\pi(y)<\infty.
\]
Let $T_x^k$ be the time of the $k$th return to $x$.
\begin{enumerate}[label=(\roman*)]
\item Show that
\[
  V_k^f=f(X(T_x^k))+\cdots+f(X(T_x^{k+1}-1)),
  \qquad k\ge1,
\]
are i.i.d. and have $E|V_k^f|<\infty$.
\item Let $K_n=\inf\{k:T_x^k\ge n\}$ and show that
\[
  \frac1n\sum_{m=1}^{K_n}V_m^f
  \longrightarrow
  \frac{EV_1^f}{E_xT_x^1}
  =
  \sum_y f(y)\pi(y)
  \qquad P_\mu\text{-a.s.}
\]
\item Show that
\[
  \frac{\max_{1\le m\le n}V_m^{|f|}}{n}\to0
\]
and conclude that
\[
  \frac1n\sum_{m=1}^n f(X_m)
  \longrightarrow
  \sum_y f(y)\pi(y)
  \qquad P_\mu\text{-a.s.}
\]
for any initial distribution $\mu$.
\end{enumerate}

\Knowledge{
\path{durrett_ch5_additive_functional_lln};
\path{durrett_ch5_stationary_measure_from_returns};
\path{durrett_ch5_positive_recurrence_return_time};
\path{durrett_ch5_strong_markov_property}.
}

\noindent\textbf{Solution.}
The return times to $x$ split the path into regenerative cycles.  By the
strong Markov property, the cycle rewards
\[
  V_k^f=\sum_{r=T_x^k}^{T_x^{k+1}-1} f(X_r),
  \qquad k\ge1,
\]
are i.i.d.

The cycle occupation measure satisfies
\[
  \mu_x(y)
  =
  E_x\sum_{r=0}^{T_x^1-1}\1_{\{X_r=y\}}.
\]
Since $\pi$ is stationary and the chain is irreducible positive
recurrent,
\[
  \mu_x(y)=\frac{\pi(y)}{\pi(x)},
  \qquad
  E_xT_x^1=\frac1{\pi(x)}.
\]
Therefore
\[
  E|V_1^f|
  \le
  \sum_y |f(y)|\mu_x(y)
  =
  \frac1{\pi(x)}\sum_y |f(y)|\pi(y)
  <\infty,
\]
and
\[
  EV_1^f
  =
  \frac1{\pi(x)}\sum_y f(y)\pi(y).
\]

By the strong law of large numbers applied to the iid cycle lengths and
cycle rewards,
\[
  \frac{T_x^k}{k}\to E_xT_x^1,
  \qquad
  \frac1k\sum_{m=1}^k V_m^f\to EV_1^f.
\]
Since $K_n/n\to 1/E_xT_x^1$, random indexing gives
\[
  \frac1n\sum_{m=1}^{K_n}V_m^f
  =
  \frac{K_n}{n}
  \cdot
  \frac1{K_n}\sum_{m=1}^{K_n}V_m^f
  \longrightarrow
  \frac{EV_1^f}{E_xT_x^1}
  =
  \sum_y f(y)\pi(y).
\]

If $Z_m=V_m^{|f|}$, then the $Z_m$ are iid and integrable.  A standard
dyadic Borel-Cantelli argument for iid integrable variables gives
\[
  \frac{\max_{1\le m\le n}Z_m}{n}\to0
  \qquad\text{a.s.}
\]
The difference between the ordinary partial sum
$\sum_{m=1}^n f(X_m)$ and the completed-cycle sum is bounded in absolute
value by the reward in at most one incomplete cycle, plus an initial
finite entrance segment if the chain does not start at $x$.  Dividing by
$n$, these errors vanish a.s.  Hence
\[
  \frac1n\sum_{m=1}^n f(X_m)
  \to
  \sum_y f(y)\pi(y)
  \qquad P_\mu\text{-a.s.}
\]

\section*{Exercise 5.6.6}

\noindent\textbf{Question.}
Central limit theorem for additive functionals.  Suppose in addition to
the conditions in Exercise 5.6.5 that
\[
  \sum_y f(y)\pi(y)=0,
  \qquad
  E_x(V_k^{|f|})^2<\infty.
\]
\begin{enumerate}[label=(\roman*)]
\item Use the random-index central limit theorem to conclude that for
any initial distribution $\mu$,
\[
  \frac1{\sqrt n}\sum_{m=1}^{K_n}V_m^f
  \Rightarrow c\chi
  \qquad\text{under }P_\mu,
\]
where $\chi$ is standard normal.
\item Show that
\[
  \frac{\max_{1\le m\le n}V_m^{|f|}}{\sqrt n}\to0
  \quad\text{in probability}
\]
and conclude that
\[
  \frac1{\sqrt n}\sum_{m=1}^n f(X_m)
  \Rightarrow c\chi
  \qquad\text{under }P_\mu.
\]
\end{enumerate}

\Knowledge{
\path{durrett_ch5_additive_functional_lln};
\path{durrett_ch3_converging_together_slutsky};
\path{durrett_ch3_lindeberg_feller_clt}.
}

\noindent\textbf{Solution.}
From Exercise 5.6.5, the cycle rewards $V_m^f$ are iid.  The assumption
$\sum_y f(y)\pi(y)=0$ implies
\[
  EV_1^f=0.
\]
The square-integrability assumption gives
\[
  \sigma^2=\operatorname{Var}(V_1^f)<\infty.
\]
Also $K_n/n\to 1/E_xT_x^1$ in probability.  The random-index central
limit theorem therefore gives
\[
  \frac1{\sqrt n}\sum_{m=1}^{K_n}V_m^f
  \Rightarrow
  N\!\left(0,\frac{\sigma^2}{E_xT_x^1}\right).
\]
Thus $c=\sigma/\sqrt{E_xT_x^1}$.

Let $Z_m=V_m^{|f|}$.  Since $EZ_1^2<\infty$,
\[
  P\left(\max_{1\le m\le n} Z_m>\varepsilon\sqrt n\right)
  \le
  nP(Z_1>\varepsilon\sqrt n).
\]
The right side tends to $0$ by the usual tail consequence of
square-integrability; equivalently, split into dyadic blocks and use
Borel-Cantelli.  Hence
\[
  \frac{\max_{1\le m\le n}V_m^{|f|}}{\sqrt n}\to0
  \quad\text{in probability}.
\]
The difference between the completed-cycle reward and the ordinary
partial sum is bounded by the reward in one incomplete cycle, plus a
finite initial segment when the initial distribution is not concentrated
at $x$.  Dividing by $\sqrt n$, this difference goes to $0$ in
probability.  Slutsky's theorem gives
\[
  \frac1{\sqrt n}\sum_{m=1}^n f(X_m)
  \Rightarrow c\chi.
\]

\section*{Exercise 5.6.7}

\noindent\textbf{Question.}
Ratio limit theorems.  Theorem 5.6.1 does not say much in the null
recurrent case.  To get a more informative limit theorem, suppose that
$y$ is recurrent and $m$ is the stationary measure on
$C_y=\{z:\rho_{yz}>0\}$, unique up to constant multiples.  Let
\[
  N_n(z)=|\{m\le n:X_m=z\}|.
\]
Break up the path at successive returns to $y$ and show that
\[
  \frac{N_n(z)}{N_n(y)}
  \longrightarrow
  \frac{m(z)}{m(y)}
  \qquad P_x\text{-a.s.}
\]
for all $x,z\in C_y$.

\Knowledge{
\path{durrett_ch5_renewal_limit_return_probabilities};
\path{durrett_ch5_stationary_measure_from_returns};
\path{durrett_ch5_stationary_measure_uniqueness_recurrent}.
}

\noindent\textbf{Solution.}
Normalize the stationary measure so that $m(y)=1$.  With this
normalization, $m=\mu_y$, the cycle occupation measure generated by
successive returns to $y$.

Let $R_k$ be the $k$th return time to $y$.  In the $k$th cycle, define
\[
  Z_k(z)=\sum_{r=R_k}^{R_{k+1}-1}\1_{\{X_r=z\}}.
\]
By the strong Markov property, the $Z_k(z)$ are iid for $k\ge1$, and
\[
  EZ_1(z)=\mu_y(z)=m(z).
\]
Also each complete cycle contains exactly one visit to $y$, namely at
its start.

By the strong law of large numbers,
\[
  \frac1K\sum_{k=1}^K Z_k(z)\to m(z).
\]
Since $N_n(y)\to\infty$ a.s. on the recurrent class and the unfinished
last cycle contributes only a negligible error when divided by
$N_n(y)$, we get
\[
  \frac{N_n(z)}{N_n(y)}\to m(z).
\]
Restoring the original normalization gives
\[
  \frac{N_n(z)}{N_n(y)}\to\frac{m(z)}{m(y)}.
\]
If the chain starts from $x\in C_y$ rather than from $y$, the finite
initial segment before the first hit of $y$ is negligible in the same
ratio.

\section*{Exercise 5.6.8}

\noindent\textbf{Question.}
We got (5.6.1) from Theorem 5.6.1 by taking expected values.  This does
not work for the ratio in the previous exercise, so we need another
approach.  Suppose $z\ne y$.
\begin{enumerate}[label=(\roman*)]
\item Let
\[
  \bar p_n(x,z)=P_x(X_n=z,T_y>n)
\]
and decompose $p^m(x,z)$ according to the value of
\[
  J=\sup\{j\in[1,m):X_j=y\}
\]
to get
\[
  \sum_{m=1}^n p^m(x,z)
  =
  \sum_{m=1}^n \bar p_m(x,z)
  +
  \sum_{j=1}^{n-1}p^j(x,y)
  \sum_{k=1}^{n-j}\bar p_k(y,z).
\]
\item Show that
\[
  \frac{\sum_{m=1}^n p^m(x,z)}
       {\sum_{m=1}^n p^m(x,y)}
  \longrightarrow
  \frac{m(z)}{m(y)}.
\]
\end{enumerate}

\Knowledge{
\path{durrett_ch5_renewal_limit_return_probabilities};
\path{durrett_ch5_hitting_return_decomposition};
\path{durrett_ch5_stationary_measure_from_returns}.
}

\noindent\textbf{Solution.}
For part (i), either the path has no visit to $y$ before time $m$, in
which case its contribution to $P_x(X_m=z)$ is $\bar p_m(x,z)$, or it
has a last visit to $y$ before time $m$.  If this last visit occurs at
time $j$, then the segment after time $j$ starts from $y$, ends at $z$
after $k=m-j$ steps, and avoids $y$ in between.  By the Markov property,
that contribution is
\[
  p^j(x,y)\bar p_k(y,z).
\]
Summing over $m\le n$ and writing $k=m-j$ gives the displayed formula.

For part (ii), normalize $m$ so that $m(y)=1$.  Then the cycle formula
for the stationary measure gives
\[
  \sum_{k=1}^\infty \bar p_k(y,z)=\mu_y(z)=m(z),
\]
because $z\ne y$.  Hence
\[
  B_N(z):=\sum_{k=1}^N \bar p_k(y,z)\longrightarrow m(z).
\]
The first term in the decomposition,
\[
  \sum_{m=1}^n\bar p_m(x,z),
\]
is bounded as $n\to\infty$ by the expected number of visits to $z$
before hitting $y$.  Meanwhile
\[
  \sum_{m=1}^n p^m(x,y)\to\infty
\]
because $y$ is recurrent and $x\in C_y$.

Using the decomposition from part (i),
\[
  \sum_{m=1}^n p^m(x,z)
  =
  O(1)+
  \sum_{j=1}^{n-1}p^j(x,y)B_{n-j}(z).
\]
This is a weighted average of $B_{n-j}(z)$ with nonnegative weights
$p^j(x,y)$, plus a negligible bounded term.  Since
$B_N(z)\to m(z)$ and the total weight
$\sum_{j=1}^{n-1}p^j(x,y)$ diverges, the Toeplitz averaging lemma gives
\[
  \frac{\sum_{m=1}^n p^m(x,z)}
       {\sum_{m=1}^n p^m(x,y)}
  \longrightarrow
  m(z).
\]
Undoing the normalization $m(y)=1$ yields
\[
  \frac{\sum_{m=1}^n p^m(x,z)}
       {\sum_{m=1}^n p^m(x,y)}
  \longrightarrow
  \frac{m(z)}{m(y)}.
\]

\section*{Exercise 5.8.1}

\noindent\textbf{Question.}
Show that $\bar X_n$ is recurrent if and only if
\[
  \sum_{n=1}^\infty \bar p^n(\alpha,\alpha)=\infty.
\]

\Knowledge{
\path{durrett_ch5_harris_recurrence};
\path{durrett_ch5_recurrence_visit_criterion};
\path{durrett_ch5_small_set_split_chain}.
}

\noindent\textbf{Solution.}
For the split chain $\bar X_n$, the artificial point $\alpha$ is an
atom.  Let
\[
  R=\inf\{n\ge1:\bar X_n=\alpha\}.
\]
By definition, the Harris chain is recurrent when
$P_\alpha(R<\infty)=1$.  Starting from $\alpha$, successive returns to
$\alpha$ regenerate the process.  Hence, exactly as in the
countable-state return criterion,
\[
  E_\alpha\sum_{n=1}^\infty \1_{\{\bar X_n=\alpha\}}
  =
  \sum_{n=1}^\infty \bar p^n(\alpha,\alpha)
  =
  \begin{cases}
    \infty, & P_\alpha(R<\infty)=1,\\
    \dfrac{P_\alpha(R<\infty)}{1-P_\alpha(R<\infty)},&
    P_\alpha(R<\infty)<1.
  \end{cases}
\]
Thus $\bar X_n$ is recurrent if and only if the displayed series
diverges.

\section*{Exercise 5.8.2}

\noindent\textbf{Question.}
If $X_n$ is a recurrent Harris chain on a countable state space, then
$S$ can only have one irreducible set of recurrent states but may have a
nonempty set of transient states.  For a concrete example, consider a
branching process in which the probability of no children $p_0>0$ and
set $A=B=\{0\}$.

\Knowledge{
\path{durrett_ch5_harris_recurrence};
\path{durrett_ch5_decomposition_theorem};
\path{durrett_ch5_branching_process_chain}.
}

\noindent\textbf{Solution.}
For a recurrent Harris chain, the split-chain atom $\alpha$ is recurrent.
Theorem 5.8.8 says that every set with positive $\lambda$-measure is
visited infinitely often from $\alpha$.  In a countable state space,
states in recurrent classes that are visible from the Harris set must
therefore communicate through the same recurrent split-chain structure.
Consequently there can be at most one irreducible recurrent class.
States outside that class may still be transient, because Harris
recurrence only forces repeated visits to sets seen by the recurrent
atom, not recurrence of every starting state.

For the branching-process example, state $0$ is absorbing.  If
$p_0>0$, then from any state $k\ge1$ there is positive probability
$p_0^k$ that all individuals have no offspring in the next generation,
so the chain can hit $0$.  Taking $A=B=\{0\}$ gives the Harris
minorization, and the split-chain atom is recurrent because once the
process reaches $0$ it stays there.

The unique recurrent class is $\{0\}$.  The states $k\ge1$ are not in
that recurrent class; when extinction occurs they are never visited
again.  Thus the recurrent Harris chain may have transient states.

\section*{Exercise 5.8.3}

\noindent\textbf{Question.}
Use Exercise 5.8.1 and imitate the proof of Theorem 5.5.10 to show that
a Harris chain with a stationary distribution must be recurrent.

\Knowledge{
\path{durrett_ch5_harris_recurrence};
\path{durrett_ch5_harris_stationary_measure};
\path{durrett_ch5_stationary_measure_definition}.
}

\noindent\textbf{Solution.}
Let $\pi$ be a stationary distribution for the original Harris chain.
Form the split-chain measure
\[
  \bar\pi=\pi\bar p.
\]
By Lemma 5.8.10, $\bar\pi$ is stationary for $\bar p$, and
\[
  \bar\pi(\alpha)=\epsilon\pi(A)>0
\]
for the Harris set $A$.

Stationarity gives $\bar\pi\bar p^n=\bar\pi$ for all $n$, so
\[
  \sum_{n=1}^\infty \bar\pi(\alpha)=\infty.
\]
Equivalently,
\[
  \int \bar\pi(dx)\sum_{n=1}^\infty \bar p^n(x,\alpha)=\infty.
\]
If $\alpha$ were transient, then the expected number of visits to
$\alpha$ starting from any state would be bounded by a finite multiple
of the probability of ever hitting $\alpha$, because after each hit the
future returns are governed by a geometric number of trials with success
probability $P_\alpha(R<\infty)<1$.  The integral above would then be
finite, a contradiction.  Hence
\[
  \sum_{n=1}^\infty \bar p^n(\alpha,\alpha)=\infty.
\]
By Exercise 5.8.1, the Harris chain is recurrent.

\section*{Exercise 5.8.4}

\noindent\textbf{Question.}
Suppose $X_n$ is a recurrent Harris chain.  Show that if $(A',B')$ is
another pair satisfying the conditions of the definition, then Theorem
5.8.8 implies
\[
  P_\alpha(\bar X_n\in A'\ \text{i.o.})=1,
\]
so the recurrence or transience does not depend on the choice of
$(A,B)$.

\Knowledge{
\path{durrett_ch5_harris_recurrence};
\path{durrett_ch5_small_set_split_chain};
\path{durrett_ch5_harris_stationary_measure}.
}

\noindent\textbf{Solution.}
Let $\lambda$ be the irreducibility measure associated with the split
chain built from $(A,B)$:
\[
  \lambda(C)=\sum_{n=1}^\infty 2^{-n}\bar p^n(\alpha,C).
\]
Because $(A',B')$ also satisfies the Harris definition, every state in
$S$ has positive probability of eventually hitting $A'$.  Starting from
$\alpha$, the split chain first redistributes into $S$ according to the
measure $\rho$ associated with $(A,B)$, and then follows the original
chain.  Therefore
\[
  P_\alpha(\tau_{A'}<\infty)>0.
\]
This implies $\lambda(A')>0$.

Since the chain is recurrent, Theorem 5.8.8 applies to every set of
positive $\lambda$-measure.  Hence
\[
  P_\alpha(\bar X_n\in A'\ \text{i.o.})=1.
\]
Thus the atom built from $(A,B)$ visits the other Harris set $A'$
infinitely often.  Reversing the roles of the two pairs gives the same
conclusion in the other split chain.  Consequently the recurrence or
transience classification is intrinsic and does not depend on the
particular choice of $(A,B)$.

\section*{Exercise 5.8.5}

\noindent\textbf{Question.}
In the GI/G/1 queue, the waiting time $W_n$ and the random walk
\[
  S_n=X_0+\xi_1+\cdots+\xi_n
\]
agree until
\[
  N=\inf\{n:S_n<0\},
\]
and at this time $W_N=0$.  Use this observation as in Example 5.3.7 to
show that Example 5.8.3 is transient when $E\xi_n>0$, recurrent when
$E\xi_n\le0$, and has a stationary distribution when $E\xi_n<0$.

\Knowledge{
\path{durrett_ch5_gig1_queue_storage_model};
\path{durrett_ch5_harris_recurrence};
\path{durrett_ch5_foster_supermartingale_recurrence};
\path{durrett_ch5_harris_stationary_measure}.
}

\noindent\textbf{Solution.}
The recursion is
\[
  W_{n+1}=(W_n+\xi_{n+1})^+.
\]
Until the ordinary random walk
\[
  S_n=X_0+\xi_1+\cdots+\xi_n
\]
first goes below $0$, the reflection has not acted, so $W_n=S_n$.  At
the first time $N=\inf\{n:S_n<0\}$, the reflected process is reset to
$0$.

If $E\xi_1>0$, then $S_n\to\infty$ with positive probability when
$X_0$ is large enough.  On that event the random walk never drops below
$0$, so the queue never returns to $0$.  Hence the Harris chain is
transient.

If $E\xi_1\le0$, the one-dimensional random walk is oscillatory or
drifts to $-\infty$, and therefore it eventually drops below $0$ a.s.
from every finite initial workload.  Thus $W_n$ hits $0$ a.s.  With
$A=B=\{0\}$ as in Example 5.8.3, this gives recurrence of the Harris
chain.

If $E\xi_1<0$, the drift toward $0$ is strict.  The same supermartingale
argument used for the M/G/1 queue gives
\[
  E_x N<\infty
\]
for every $x$.  Thus the return time to the Harris atom has finite mean,
so the recurrent Harris chain is positive recurrent and has a stationary
distribution.

\section*{Exercise 5.8.6}

\noindent\textbf{Question.}
Let
\[
  m_n=\min(S_0,S_1,\ldots,S_n),
\]
where $S_n$ is the random walk defined previously.
\begin{enumerate}[label=(\roman*)]
\item Show that
\[
  S_n-m_n\stackrel{d}{=}W_n.
\]
\item Let $\xi'_m=\xi_{n+1-m}$ for $1\le m\le n$.  Show that
\[
  S_n-m_n=\max(S'_0,S'_1,\ldots,S'_n).
\]
\item Conclude that as $n\to\infty$,
\[
  W_n\Rightarrow M\equiv\max(S'_0,S'_1,S'_2,\ldots).
\]
\end{enumerate}

\Knowledge{
\path{durrett_ch5_gig1_queue_storage_model};
\path{durrett_ch5_random_walk_chain};
\path{durrett_ch3_weak_convergence_definition}.
}

\noindent\textbf{Solution.}
Assume $W_0=S_0=0$.  The reflected recursion has the explicit
Skorokhod-reflection form
\[
  W_n=S_n-\min_{0\le j\le n}S_j
\]
as can be checked by induction.  Indeed, if the running minimum does not
decrease at time $n+1$, then both sides add $\xi_{n+1}$; if the running
minimum decreases, then the right side becomes $0$, matching the
positive-part recursion.  Hence
\[
  W_n=S_n-m_n.
\]
This proves the distributional statement, and in this construction the
two sides are actually equal.

For part (ii),
\[
  S_n-m_n
  =
  \max_{0\le j\le n}(S_n-S_j)
  =
  \max_{0\le j\le n}(\xi_{j+1}+\cdots+\xi_n).
\]
If $\xi'_m=\xi_{n+1-m}$ and
$S'_k=\xi'_1+\cdots+\xi'_k$, then the sums
$\xi_{j+1}+\cdots+\xi_n$ are exactly
$S'_{n-j}$ for $0\le j\le n$.  Therefore
\[
  S_n-m_n=\max_{0\le k\le n}S'_k.
\]

If $E\xi_1<0$, then the reversed increments have the same distribution
as the original increments and the random walk $S'_k$ drifts to
$-\infty$.  Hence
\[
  M=\max_{k\ge0}S'_k<\infty
  \qquad\text{a.s.}
\]
The finite maxima $\max_{0\le k\le n}S'_k$ increase a.s. to $M$.
Since $W_n$ has the same distribution as this finite maximum,
\[
  W_n\Rightarrow M.
\]

\section*{Exercise 5.8.7}

\noindent\textbf{Question.}
In the discrete O.U. process, $X_{n+1}$ is normal with mean
$\theta X_n$ and variance $1$.  What happens to the recurrence and
transience if instead $Y_{n+1}$ is normal with mean $0$ and variance
$\beta^2|Y_n|$?

\Knowledge{
\path{durrett_ch5_harris_recurrence};
\path{durrett_ch5_harris_convergence_theorem};
\path{durrett_ch5_foster_supermartingale_recurrence}.
}

\noindent\textbf{Solution.}
If $Y_n=0$, then the conditional variance is $0$, so $0$ is absorbing.
Starting from $Y_0\ne0$, the transition distribution is continuous and
puts probability zero on the single point $0$.  Thus $\{0\}$ is a
separate absorbing recurrent class, while the nonzero states form the
interesting continuous class.

For $Y_n\ne0$, write $R_n=|Y_n|$ and $L_n=\log R_n$.  If $Z_{n+1}$ is a
standard normal random variable, then
\[
  Y_{n+1}=\beta\sqrt{|Y_n|}\,Z_{n+1},
\]
so
\[
  L_{n+1}
  =
  \log\beta+\frac12L_n+\log|Z_{n+1}|.
\]
This is an autoregressive recursion on the log-scale with coefficient
$1/2$.  Since $E|\log|Z||<\infty$, the recursion is stable and has a
proper stationary distribution:
\[
  L_\infty
  \stackrel{d}{=}
  \sum_{j=0}^\infty 2^{-j}
  \bigl(\log\beta+\log|Z_{j+1}|\bigr).
\]
Thus $|Y_n|$ neither escapes to infinity nor collapses to zero from a
nonzero start; it returns repeatedly to bounded log-scale sets.

Consequently, for every $\beta>0$, the chain restricted to
$\mathbb{R}\setminus\{0\}$ is recurrent in the Harris sense and in fact
has a stationary distribution.  The full chain is not irreducible
because $0$ is absorbing and cannot be reached from nonzero states, but
its recurrence structure is clear: $\{0\}$ is an absorbing recurrent
class, and the nonzero Harris class is positive recurrent.

\end{document}
